\chapter{Introduction: The API Economy \& How to Use This Book}
\label{ch:introduction}

\section*{Executive Summary}
\addcontentsline{toc}{section}{Executive Summary}

Every system you will build, operate, or integrate over the next decade will
either expose an Application Programming Interface (API) or depend on one ---
usually both. The modern enterprise is no longer a monolith behind a firewall;
it is a mesh of services communicating over versioned contracts, where the
interface \emph{is} the product and the contract \emph{is} the engineering
artifact of record. Industry analysts have estimated for years that the large
majority of enterprise workloads now traverse an API boundary at some point in
their lifecycle, and the economic consequences are measurable: entire
companies (payments processors, communications platforms, mapping providers)
sell nothing except well-designed interfaces over HTTP. This is the
\emph{API economy}: interfaces priced, versioned, secured, and observed with
the same rigor once reserved for physical supply chains.

This book is an engineering curriculum for that reality. It treats APIs not
as an afterthought bolted onto an application but as \textbf{products},
\textbf{contracts}, and \textbf{failure surfaces} --- three mental models
that recur in every chapter. As products, APIs have consumers, developer
experience, pricing, and lifecycles. As contracts, they have formal,
machine-checkable semantics: operations, representations, status codes, and
failure behavior, all versioned over time. As failure surfaces, they are the
exact location where distributed systems hurt you: timeouts, retries, partial
failures, and misaligned assumptions between provider and consumer. Mastering
APIs means mastering all three perspectives simultaneously.

You are reading Chapter 0, the orientation chapter. Its job is fourfold.
First, it defines terms \emph{formally} --- we will not wave our hands at
words like ``REST,'' ``contract,'' or ``latency percentile.'' Second, it maps
the territory: four modules, twenty-one technical chapters, and the
dependency structure between them, so you can plan a path whether you are a
consumer of APIs, a producer of them, an architect governing a platform, or a
data engineer moving records across boundaries. Third, it installs the
evaluation discipline used throughout the book: latency measured in
percentiles rather than averages, error budgets derived from service-level
objectives, and benchmark protocols that a skeptical reviewer can reproduce.
Fourth, it gets your hands dirty within the first three hours: an offline,
stdlib-only lab in which you build a working HTTP API and client on your own
machine and measure its latency distribution.

\begin{keyconceptbox}
An API is a \textbf{versioned contract} between a provider and its consumers.
The contract specifies \emph{operations} (what you may ask), \emph{representations}
(the shape of data exchanged), and \emph{failure semantics} (what happens when
things go wrong), all carried over a \emph{transport protocol}. Everything else
in this book --- frameworks, gateways, meshes, observability --- exists to make
that contract reliable, secure, and evolvable at scale.
\end{keyconceptbox}

The tone throughout is pragmatic and production-focused. Every technique is
introduced because it prevents a real failure mode, and every chapter closes
with measurable success criteria rather than vague encouragement. Code is
Python 3.11+, fully typed, stdlib-first, and runnable offline --- no cloud
accounts, no API keys, no network access required. Where third-party services
appear (payment gateways, identity providers), they are simulated with
deterministic local stubs. All datasets are synthetic, drawn from the
fictional \emph{Northwind Robotics} universe used across this curriculum.

By the end of this chapter you will have a precise vocabulary, a working local
environment, a verified benchmark harness, and a map of the road ahead. Let us
begin.

%% ============================================================================
\section{Learning Objectives \& Prerequisites}
\label{sec:ch0-objectives}

\subsection{Learning Objectives}

After completing this chapter and its lab, you will be able to:

\begin{enumerate}[label=\textbf{LO-\arabic*.}]
  \item Define an API formally as a versioned contract (operations,
        representations, failure semantics, transport) and explain the API
        value chain from provider to consumer to end user.
  \item Apply the three mental models --- API as product, API as contract,
        API as failure surface --- to reason about design decisions.
  \item Classify real incidents into the five-category API failure taxonomy
        (transport, contract, semantic, operational, security) and select an
        appropriate detection strategy for each.
  \item Locate any system on the five-level API Maturity Model (L0--L4) and
        articulate the concrete capabilities required to advance one level.
  \item Measure client-observed latency correctly using percentiles
        (p50/p95/p99), explain why arithmetic means mislead, and run the
        minimal benchmark protocol defined in this chapter.
  \item Decompose a request's latency budget across TLS setup, validation,
        business logic, and serialization, and compute an error budget from
        an SLO target (e.g.\ 99.9\% monthly availability $= 43.2$ minutes of
        downtime).
  \item Set up a reproducible Python development environment on Windows
        PowerShell and verify it with an automated sanity-check script.
  \item Build, run, and measure a minimal stdlib-only HTTP API and client
        offline, forming the baseline for every subsequent chapter.
  \item Choose the correct reader pathway (consumer, producer, architect,
        data engineer) and locate prerequisite chapters for any topic.
  \item State the governance, threat-modeling (STRIDE), and ethics baselines
        that frame all production work in this book.
\end{enumerate}

\subsection{Prerequisites}

This is an intensive curriculum aimed at practicing engineers. We assume:

\begin{itemize}
  \item \textbf{Python 3.11 or newer.} You should be comfortable with
        functions, classes, type hints (PEP 484), context managers, and
        virtual environments. You do \emph{not} need prior web-framework
        experience; FastAPI is taught from first principles in
        Chapter~\ref{ch:fastapi}.
  \item \textbf{Basic networking literacy.} You should know what an IP
        address, a TCP connection, a DNS lookup, and a port are. We teach
        HTTP from the wire up in Chapter~\ref{ch:http_wire}, so deep protocol
        knowledge is \emph{not} a prerequisite --- but nothing in that
        chapter should be your first exposure to the word ``packet.''
  \item \textbf{Git.} Cloning, branching, committing. All sample code is
        versioned, and the labs assume you can diff and revert.
  \item \textbf{A Windows workstation with PowerShell.} All setup and run
        instructions are given in PowerShell idioms. Everything also works on
        macOS/Linux, but those commands are intentionally out of scope.
\end{itemize}

\subsection{How to Use This Textbook}

The book supports three modes of use, and you may move between them freely:

\begin{description}
  \item[Self-paced study (16 weeks).] One module per four weeks, one to two
    chapters per week. Do every lab; the labs are where the learning
    consolidates. Budget 6--10 hours per week including exercises.
  \item[Team workshops.] Each chapter's lab and capstone are scoped for a
    half-day workshop. A common pattern is a ``guild'' format: one engineer
    presents the chapter's failure taxonomy section, the team then runs the
    lab against a shared repository, and the session closes with the
    self-assessment questions answered collaboratively.
  \item[Reference.] Chapters are cross-referenced rather than strictly
    linear. If you arrive with production experience, jump directly to the
    chapter you need and use the dependency map
    (Section~\ref{sec:ch0-depmap}) to backfill prerequisites.
\end{description}

\subsection{Reader Pathways}
\label{sec:ch0-pathways}

Four curated tracks through the same twenty-one chapters. Every track reads
Chapters 0--4; divergence is deliberate, and each track rejoins the shared
core where its role demands it.

\begin{table}[h]
  \centering
  \small
  \begin{tabularx}{\textwidth}{l X}
    \toprule
    \textbf{Track} & \textbf{Chapter sequence} \\
    \midrule
    Consumer-only &
      0--5, 8 (as a client), 10, 14 (as a client), 19 (consumer contract
      tests). Skip provider build chapters; read their ``contract'' sections
      only. \\
    Producer &
      0--10, 16--20. Add 11--13 selectively per your product's protocol
      surface. This is the default full-stack track. \\
    Architect &
      0--1, 3--4, 9--15, 17--21. Skim implementation detail in 2 and 6--8;
      go deep on versioning, resilience, observability, and governance. \\
    Data engineer &
      0--5, 8, 10--11, 14--15, 18. Emphasis on pagination, idempotent
      pipelines, webhooks/streams, and backpressure; GraphQL (12) if your
      warehouse feeds BI consumers. \\
    \bottomrule
  \end{tabularx}
  \caption{Reader pathways. Chapter numbers refer to the dependency map in
    Table~\ref{tab:ch0-deps}.}
  \label{tab:ch0-pathways}
\end{table}

\subsection{Conventions Used in This Book}

\begin{itemize}
  \item \textbf{Code listings} are complete and runnable. Nothing is elided;
    there are no ``left as an exercise'' gaps in production listings.
    Python listings are typed (PEP 484) and lint-clean (PEP 8).
  \item \textbf{Boxes.} \emph{Key Concept} boxes state load-bearing ideas you
    must internalize. \emph{Warning} boxes mark mistakes that cause
    production incidents. \emph{Example} boxes ground abstractions in the
    Northwind Robotics universe.
  \item \textbf{Formal definitions} appear in numbered Definition blocks and
    are the canonical meaning of the term for the rest of the book.
  \item \textbf{Math notation.} Latency percentiles are written
    $p_{50}, p_{95}, p_{99}$; rates are written as $\lambda$ (arrival rate,
    requests/second) and $\mu$ (service rate); queueing formulas are stated
    precisely when used. Scalars are italic ($x$), vectors bold
    ($\mathbf{x}$).
  \item \textbf{Paths and commands} are Windows-native. In prose, a path such
    as the virtual-environment activation script is written
    \texttt{.venv\textbackslash Scripts\textbackslash Activate.ps1}; inside
    code blocks, backslashes are literal.
  \item \textbf{Citations} refer to the bibliography; foundational references
    such as Fielding's dissertation \cite{fielding2000rest} and the HTTP
    semantics standard \cite{rfc9110} are cited at first substantive use.
\end{itemize}

\begin{warningbox}
\textbf{Offline-first is a hard constraint, not a suggestion.} Every listing,
lab, and test in this book runs with the network cable unplugged. If you find
yourself reaching for a live third-party API while doing an exercise, stop:
the chapter provides a deterministic local stub for it. This discipline is
deliberate --- it forces you to understand the protocol rather than the
vendor SDK, and it makes every result in the book reproducible on a fresh
machine.
\end{warningbox}
%% ============================================================================
\section{Deep Technical Mechanics \& Architectural Concepts}
\label{sec:ch0-mechanics}

\subsection{What an API Is: A Formal Definition}

The word ``API'' is used so loosely in industry that it has nearly lost its
meaning. In this book it has exactly one meaning, and we will use it
consistently for twenty-one chapters.

\begin{definition}[Application Programming Interface]
\label{def:api}
An \emph{API} is a \textbf{versioned contract} between a \emph{provider} and
one or more \emph{consumers}, specifying:
\begin{enumerate}
  \item a set of \textbf{operations} $O = \{o_1, o_2, \dots, o_n\}$, each with
        a name, input parameters, and preconditions;
  \item \textbf{representations} $R$: the serialized data shapes exchanged in
        requests and responses (schemas, media types, encodings);
  \item \textbf{failure semantics} $F$: for each operation, the complete set
        of failure modes (status codes, error representations, retryability,
        and side-effect guarantees); and
  \item a \textbf{transport protocol} $T$ over which operations are invoked
        (e.g.\ HTTP/1.1 or HTTP/2 per RFC~9110 \cite{rfc9110}, gRPC over
        HTTP/2, or a message broker binding).
\end{enumerate}
A \emph{version} $v$ of the contract is an immutable snapshot
$\langle O_v, R_v, F_v, T \rangle$; evolving any component produces a new
version governed by a compatibility policy (Chapter~\ref{ch:versioning}).
\end{definition}

Three consequences of Definition~\ref{def:api} are worth stating now because
they drive everything that follows. First, an API is \emph{not} the
implementation; two services with identical contracts are substitutable,
which is precisely what makes mocking, contract testing, and the
offline-first discipline of this book possible. Second, failure semantics are
part of the contract --- a provider that changes which errors it returns has
broken its contract just as surely as one that renames a field. Third,
versioning is not optional metadata; it is part of the definition. An
unversioned API is an incomplete contract whose terms can change under the
consumer at any moment.

\subsection{The API Value Chain}

APIs exist to move value across organizational and machine boundaries. The
canonical value chain has five links:

\begin{enumerate}
  \item \textbf{Capability.} Something the provider can do: charge a card,
        forecast demand, retrieve a robot's telemetry.
  \item \textbf{Interface.} The versioned contract exposing the capability.
  \item \textbf{Distribution.} How consumers discover and integrate:
        documentation, catalogs, SDKs, developer portals.
  \item \textbf{Consumption.} Client applications composed from one or more
        APIs.
  \item \textbf{End-user value.} The business outcome, measured in revenue,
        cost avoided, or risk reduced.
\end{enumerate}

Value leaks at every seam. A capability without a contract is unreachable; a
contract without distribution is undiscovered; consumption without
observability is unbillable and un-debuggable. The maturity model in
Section~\ref{sec:ch0-maturity} is, at heart, a checklist for sealing those
seams. Treating the API \emph{as a product} \cite{gehl2021apidesign} means
assigning an owner to the whole chain, not merely to link two.

\begin{examplebox}
\textbf{Northwind Robotics.} Northwind builds warehouse robots. Its telemetry
platform (capability) exposes \texttt{GET /v1/robots/\{id\}/telemetry}
(interface), published through an internal developer portal with an OpenAPI
3.1 document \cite{openapi31} (distribution), consumed by the fleet-dashboard
web app and a nightly data-pipeline job (consumption), which together cut
unplanned robot downtime by 18\% last quarter (end-user value). When the
pipeline team complains that ``the API is slow,'' every link of this chain is
a suspect --- and this chapter's diagnostic playbook
(Section~\ref{sec:ch0-playbook}) tells you how to isolate which one.
\end{examplebox}

\subsection{Three Mental Models: Product, Contract, Failure Surface}

\begin{description}
  \item[API as product.] The API has customers (developers), a funnel
    (discovery $\to$ first successful call $\to$ production traffic),
    competitors, and a lifecycle from beta to deprecation and sunset
    (Chapter~\ref{ch:lifecycle}). Product thinking forces you to budget for
    documentation, changelogs, and developer experience (DX) as
    first-class engineering deliverables.
  \item[API as contract.] The contract is machine-checkable. OpenAPI
    documents, JSON Schemas, and protocol-buffer definitions are not
    documentation \emph{about} the system; they are artifacts \emph{of} the
    system, diffed in code review and enforced in CI
    (Chapters~\ref{ch:serialization} and \ref{ch:testing}). Contract
    thinking turns ``we changed a field, hope nobody noticed'' into a build
    failure.
  \item[API as failure surface.] Every API call is a distributed-systems
    event: it can be delayed, duplicated, dropped, reordered, or partially
    applied. The consumer and provider hold independent, possibly
    inconsistent views of the world \cite{kleppmann2017ddia}. Failure-surface
    thinking is why this book spends entire chapters on idempotency
    (Chapter~\ref{ch:idempotency}), rate limiting
    (Chapter~\ref{ch:rate_limiting}), and resilience patterns
    (Chapter~\ref{ch:microservices}).
\end{description}

\subsection{The Five-Category API Failure Taxonomy}
\label{sec:ch0-taxonomy}

Debugging is classification before it is repair. Nearly every API incident
you will ever meet belongs to exactly one of five categories. Learn the
taxonomy now; every subsequent chapter deepens one or more rows of it.

\begin{table}[h]
  \centering
  \small
  \begin{tabularx}{\textwidth}{l X X}
    \toprule
    \textbf{Category} & \textbf{Example} & \textbf{Primary detection strategy} \\
    \midrule
    Transport &
      TCP connect timeout; TLS handshake failure; DNS resolution error;
      HTTP 502 from an intermediate proxy. &
      Connection-level metrics, synthetic probes, TCP/TLS handshake
      timing, load-balancer logs. \\
    Contract &
      Provider renames \texttt{unit\_price} to \texttt{unitPrice}; consumer
      deserializer throws; HTTP 400 on previously valid payload. &
      Schema-diff gates in CI, consumer-driven contract tests, canary
      traffic replay. \\
    Semantic &
      Call succeeds (HTTP 200) but means the wrong thing: price returned in
      cents instead of dollars; ``deleted'' flag inverted. &
      Property-based tests, reconciliation jobs, business-metric anomaly
      detection. \\
    Operational &
      Latency p99 degrades from 120\,ms to 4\,s under load; retry storm
      after a partial outage; connection-pool exhaustion. &
      Percentile latency dashboards, saturation metrics, distributed
      tracing, load tests. \\
    Security &
      Broken object-level authorization (BOLA); leaked API key in a client
      bundle; unsanitized input reaching a SQL query. &
      Authz fuzzing, secret scanning, WAF/IDS alerts, penetration tests,
      OWASP API Top 10 audits. \\
    \bottomrule
  \end{tabularx}
  \caption{The API failure taxonomy used throughout this book.}
  \label{tab:ch0-taxonomy}
\end{table}

Two rules make the taxonomy actionable. Rule one: \textbf{classify before you
mitigate.} Retrying a contract failure is waste; retrying a transport failure
with exponential backoff is correct; retrying a non-idempotent semantic
operation is a corruption vector. Rule two: \textbf{each category has a
natural home in the stack.} Transport failures live below HTTP semantics,
contract failures at the schema layer, semantic failures in business logic,
operational failures in capacity and dependency management, and security
failures at every boundary --- which is why Chapter~\ref{ch:api_security}
treats the OWASP API Security Top 10 as a cross-cutting concern rather than a
bolt-on.

\subsection{The API Maturity Model (L0--L4)}
\label{sec:ch0-maturity}

Organizations do not leap from shell scripts to governed platforms; they
climb. The five-level model below is the book's yardstick --- you will be
asked to place real systems on it in the exercises.

\begin{description}
  \item[L0 --- Ad-hoc scripts.] Calls are made with one-off scripts, copied
    curl commands, and hardcoded credentials. No versioning, no tests, no
    runbooks. Characteristic quote: ``it works on my machine.''
  \item[L1 --- Documented REST.] Endpoints follow REST constraints
    \cite{fielding2000rest}, ship human-readable documentation, and use
    structured errors. Contracts exist but are enforced by convention, not
    by machines.
  \item[L2 --- Contract-tested in CI.] OpenAPI/JSON Schema artifacts are
    generated or validated in CI; breaking changes fail the build; consumer
    contract tests run before deploys (Chapter~\ref{ch:testing}).
  \item[L3 --- Resilient, observable services.] Timeouts, retries with
    budgets, circuit breakers, and idempotency keys are standard
    \cite{nygard2018releaseit}; SLOs with error budgets are defined and
    alerted on \cite{beyer2016sre}; traces and structured logs make every
    request explainable (Chapters \ref{ch:observability} and
    \ref{ch:microservices}).
  \item[L4 --- Governed API platform.] A developer portal, catalog,
    lifecycle policies, automated governance linting, security baselines,
    and paved-road templates make the compliant path the easy path
    (Chapter~\ref{ch:lifecycle}).
\end{description}

\begin{figure}[h]
  \centering
  \begin{tikzpicture}[
      font=\small,
      level/.style={draw, rounded corners, minimum width=11.8cm,
                    minimum height=1.05cm, align=center},
      arr/.style={->, >=stealth, thick, packtgray}
    ]
    \node[level, fill=red!8]     (l0) {\textbf{L0 Ad-hoc scripts} --- curl in a wiki page; secrets in code};
    \node[level, fill=orange!10, above=0.35cm of l0] (l1) {\textbf{L1 Documented REST} --- style guide, docs, structured errors};
    \node[level, fill=yellow!12, above=0.35cm of l1] (l2) {\textbf{L2 Contract-tested CI} --- schema diff gates, consumer contracts};
    \node[level, fill=green!10, above=0.35cm of l2] (l3) {\textbf{L3 Resilient \& observable} --- SLOs, error budgets, tracing, breakers};
    \node[level, fill=blue!8, above=0.35cm of l3] (l4) {\textbf{L4 Governed platform} --- portal, catalog, lifecycle policy, paved road};
    \draw[arr] (l0) -- (l1);
    \draw[arr] (l1) -- (l2);
    \draw[arr] (l2) -- (l3);
    \draw[arr] (l3) -- (l4);
    \node[right=0.25cm of l0, anchor=west] {\footnotesize Module 1};
    \node[right=0.25cm of l2, anchor=west] {\footnotesize Modules 2--3};
    \node[right=0.25cm of l4, anchor=west] {\footnotesize Module 4};
  \end{tikzpicture}
  \caption{The API Maturity Model. Each level is a prerequisite for the next;
    the right margin shows which module of this book builds that level's
    capabilities.}
  \label{fig:ch0-maturity}
\end{figure}

\subsection{Map of the Book: Four Modules and a Synthesis}
\label{sec:ch0-depmap}

The curriculum is a 16-week intensive organized into four modules plus a
synthesis chapter.

\paragraph{Module 1 --- Foundations (Chapters 1--5).}
Chapter~\ref{ch:http_wire} opens the wire: requests, responses, headers, and
status codes per RFC~9110 \cite{rfc9110}, captured and decoded by hand.
Chapter~\ref{ch:consuming_python} builds disciplined HTTP clients in Python:
sessions, timeouts, retries, and connection pooling.
Chapter~\ref{ch:rest_style} formalizes REST as an architectural style
\cite{fielding2000rest} --- resources, representations, hypermedia, and the
constraints that actually matter. Chapter~\ref{ch:serialization} covers
JSON, schema languages, and contract artifacts. Chapter~\ref{ch:client_security}
secures the consumer side: credential storage, TLS verification, and token
handling.

\paragraph{Module 2 --- Designing \& Building (Chapters 6--10).}
Chapter~\ref{ch:fastapi} builds production services with FastAPI.
Chapter~\ref{ch:problem_details} standardizes errors with RFC~9457 problem
details. Chapter~\ref{ch:pagination} solves pagination, filtering, and
sorting over large collections. Chapter~\ref{ch:versioning} is the contract
evolution chapter: compatibility classes, version negotiation, deprecation
policy. Chapter~\ref{ch:idempotency} closes the module with the trinity of
safe repetition: idempotency keys, HTTP caching, and concurrency control.

\paragraph{Module 3 --- Advanced Architecture (Chapters 11--15).}
Chapter~\ref{ch:async_apis} covers webhooks, Server-Sent Events, and
WebSockets. Chapters \ref{ch:graphql} and \ref{ch:grpc} teach the two major
alternative paradigms --- GraphQL's client-driven queries and gRPC's
schema-first, high-performance RPC. Chapter~\ref{ch:rate_limiting} derives
rate limiting and throttling algorithms from queueing first principles.
Chapter~\ref{ch:microservices} assembles resilience patterns for
service-to-service calls \cite{newman2021microservices}.

\paragraph{Module 4 --- Enterprise Production \& Security (Chapters 16--20).}
Chapter~\ref{ch:api_security} works through the OWASP API Security Top 10
with exploitable labs and fixes. Chapter~\ref{ch:deployment} covers
deployment, gateways, and Kubernetes. Chapter~\ref{ch:observability} is SRE
applied to APIs: SLIs, SLOs, error budgets, tracing \cite{beyer2016sre}.
Chapter~\ref{ch:testing} is quality engineering: contract, property-based,
chaos, and load testing. Chapter~\ref{ch:lifecycle} closes with governance,
lifecycle management, and developer experience.

\paragraph{Synthesis (Chapter 21).}
Chapter~\ref{ch:synthesis} integrates everything into a capstone
architecture and looks ahead at emerging protocol and platform trends.

\paragraph{Chapter dependency map.}
Solid knowledge of a chapter requires its ancestors. The edges below are the
\emph{hard} prerequisites; all other cross-references are enrichment.

\begin{table}[h]
  \centering
  \small
  \begin{tabularx}{\textwidth}{l l X}
    \toprule
    \textbf{Chapter} & \textbf{Requires} & \textbf{Why} \\
    \midrule
    1 HTTP wire & 0 & Baseline environment and vocabulary. \\
    2 Consuming APIs & 1 & Clients are built on raw protocol semantics. \\
    3 REST style & 1 & Constraints are stated in protocol terms. \\
    4 Serialization \& contracts & 2--3 & Schemas describe REST payloads. \\
    5 Client-side security & 2, 4 & Secures typed, contract-aware clients. \\
    6 FastAPI & 3--4 & Framework routes map to resources and schemas. \\
    7 Problem details & 6 & Error media types plug into the framework. \\
    8 Pagination/filter/sort & 4, 6 & Collection contracts over real services. \\
    9 Versioning & 4, 7 & Compatibility is defined over schemas \& errors. \\
    10 Idempotency/caching & 1, 6--7 & Safe retries need protocol \& error semantics. \\
    11 Async APIs & 1, 6 & Webhooks/SSE/WebSockets extend the HTTP mental model. \\
    12 GraphQL & 3--4 & A different contract style over the same wire. \\
    13 gRPC & 4 & Schema-first RPC; protobuf is a serialization contract. \\
    14 Rate limiting & 2, 10 & Builds on client behavior and safe retries. \\
    15 Microservices resilience & 10, 14 & Breakers and budgets compose earlier patterns. \\
    16 API security & 5, 6 & Threats are mapped onto real services and clients. \\
    17 Deployment \& gateways & 6, 15 & Ships resilient services behind gateways. \\
    18 Observability \& SRE & 14--15 & SLOs sit on rate-limited, resilient services. \\
    19 Testing \& quality & 4, 6--7 & Contract tests need schemas and error models. \\
    20 Lifecycle \& governance & 9, 16, 19 & Governance composes versioning, security, testing. \\
    21 Synthesis & all & Integrative capstone. \\
    \bottomrule
  \end{tabularx}
  \caption{Hard prerequisite edges between chapters.}
  \label{tab:ch0-deps}
\end{table}

\subsection{Reference Architecture Variants}
\label{sec:ch0-refarch}

Four deployment topologies recur across the book. Each chapter states which
variant its listings target; the evaluation protocol
(Section~\ref{sec:ch0-eval}) is defined identically for all four so that
results are comparable.

\begin{description}
  \item[V1 --- Client-only consumer.] Your code calls an external (simulated)
    provider. Used in Chapters \ref{ch:http_wire}--\ref{ch:client_security}.
  \item[V2 --- Single FastAPI service.] One process, one datastore, local
    clients. Used in Chapters \ref{ch:fastapi}--\ref{ch:idempotency}.
  \item[V3 --- Gateway-fronted microservices.] An edge gateway routes to
    multiple services with resilience policies. Used in Chapters
    \ref{ch:rate_limiting}--\ref{ch:observability}.
  \item[V4 --- Event-driven.] Producers and consumers decoupled through
    topics; webhooks and streams. Used in Chapter~\ref{ch:async_apis} and
    parts of \ref{ch:lifecycle}.
\end{description}

\begin{figure}[h]
  \centering
  \begin{tikzpicture}[
      font=\small,
      box/.style={draw, rounded corners, minimum height=0.95cm,
                  align=center, fill=blue!6},
      datastore/.style={draw, cylinder, shape border rotate=90,
                        aspect=0.28, minimum height=1.0cm, align=center,
                        fill=green!8},
      arr/.style={->, >=stealth, thick},
      lbl/.style={midway, above, font=\footnotesize, align=center}
    ]
    % Actors
    \node[box, minimum width=2.2cm] (client) {\textbf{Client}\\consumer app};
    \node[box, minimum width=2.6cm, right=1.6cm of client, fill=orange!12]
      (gateway) {\textbf{Edge gateway}\\authn, routing, quotas};
    \node[box, minimum width=2.4cm, right=1.6cm of gateway] (svcA)
      {\textbf{Service A}\\orders (FastAPI)};
    \node[box, minimum width=2.4cm, below=1.1cm of svcA] (svcB)
      {\textbf{Service B}\\telemetry (FastAPI)};
    \node[datastore, right=1.5cm of svcA] (dbA) {DB\\A};
    \node[datastore, right=1.5cm of svcB] (dbB) {DB\\B};
    \node[box, minimum width=2.6cm, below=1.1cm of svcB, fill=yellow!14]
      (bus) {\textbf{Event bus}\\topics / webhooks};
    % Edges
    \draw[arr] (client) -- node[lbl]{HTTPS / JSON} (gateway);
    \draw[arr] (gateway) -- node[lbl]{route \texttt{/orders}} (svcA);
    \draw[arr] (gateway) |- node[lbl, near end]{route \texttt{/telemetry}} (svcB);
    \draw[arr] (svcA) -- (dbA);
    \draw[arr] (svcB) -- (dbB);
    \draw[arr] (svcA) |- node[lbl, near start]{publish events} (bus);
    \draw[arr] (bus) -| node[lbl, near end, below]{push / webhook} (client);
    % Variant markers
    \node[draw, densely dashed, rounded corners, inner sep=0.12cm,
          fit=(client)(svcA), label={[font=\footnotesize]south:{V1--V2 collapse to client $\to$ one service}}] {};
  \end{tikzpicture}
  \caption{Reference architecture (variant V3 shown with V4 event leg).
    Variant V1 uses only the client against a stubbed provider; V2 removes
    the gateway and second service.}
  \label{fig:ch0-refarch}
\end{figure}

\subsection{Cost and Latency Budgeting}
\label{sec:ch0-budget}

A request's end-to-end latency $L$ decomposes into additive stages:

\begin{equation}
  L \;=\; L_{\mathrm{TLS}} + L_{\mathrm{queue}} + L_{\mathrm{validate}}
        + L_{\mathrm{logic}} + L_{\mathrm{serialize}} + L_{\mathrm{egress}}
  \label{eq:ch0-budget}
\end{equation}

where $L_{\mathrm{TLS}}$ is connection setup (amortized to near zero by
keep-alive), $L_{\mathrm{queue}}$ is time waiting for a worker,
$L_{\mathrm{validate}}$ is schema and authz checking, $L_{\mathrm{logic}}$ is
business logic including downstream calls, and
$L_{\mathrm{serialize}} + L_{\mathrm{egress}}$ is response encoding and
transmission. A \emph{latency budget} assigns each stage a target share of
the SLO; Table~\ref{tab:ch0-budget} shows the starter budget used in labs.
Cost budgeting mirrors the same decomposition: TLS and egress are dominated
by network, validation and serialization by CPU, logic by I/O and downstream
fan-out. When p99 violates the SLO,
Equation~\eqref{eq:ch0-budget} tells you where to spend optimization effort
\emph{before} you profile anything: measure each stage, compare to its
budget, and fix the largest overrun first.

\begin{table}[h]
  \centering
  \small
  \begin{tabularx}{\textwidth}{l c X}
    \toprule
    \textbf{Stage} & \textbf{Starter budget} & \textbf{Dominant cost driver} \\
    \midrule
    TLS / connection setup & $\leq 10\%$ & Handshake round trips; keep-alive hit rate. \\
    Queueing & $\leq 10\%$ & Worker pool saturation; arrival bursts. \\
    Request validation & $\leq 10\%$ & Schema complexity; CPU per request. \\
    Business logic \& I/O & $\leq 55\%$ & Downstream fan-out; query cost. \\
    Serialization \& egress & $\leq 15\%$ & Payload size; compression ratio. \\
    \bottomrule
  \end{tabularx}
  \caption{Starter latency budget decomposition (share of p99 target).}
  \label{tab:ch0-budget}
\end{table}

\subsection{Governance, SLOs, and Threat Modeling: The Baseline}
\label{sec:ch0-governance}

\paragraph{Governance baseline.}
Every API in this book, including toy ones, carries five artifacts: an owner,
a version, a machine-readable contract, a changelog, and a deprecation
policy. That is the minimum bar for L1 on the maturity ladder, and it costs
almost nothing at small scale. Governance at L4 adds catalogs, automated
linting of contracts, and lifecycle state machines (Chapter~\ref{ch:lifecycle}).

\paragraph{SLO and error-budget starter.}
A service-level objective states a target for a measurable indicator. The
canonical first SLO is availability: the fraction of requests answered with a
non-5xx status within the latency threshold. The \emph{error budget} is the
allowed shortfall. Worked example: a 99.9\% monthly availability SLO over a
30-day month permits
\begin{equation}
  (1 - 0.999) \times 30 \times 24 \times 60 \;=\; 43.2
  \ \text{minutes of downtime per month.}
  \label{eq:ch0-errorbudget}
\end{equation}
The budget converts reliability from a vibe into a spendable resource: when
more than half the budget is gone mid-month, you freeze feature releases and
spend engineering on reliability instead \cite{beyer2016sre}. We compute
budgets for every lab from Chapter~\ref{ch:observability} onward.

\paragraph{Threat model primer (STRIDE).}
Before any endpoint ships, ask six questions --- one per STRIDE category
\cite{gehl2021apidesign}:

\begin{itemize}
  \item \textbf{S}poofing: can a caller impersonate another principal?
        (Control: authentication on every request; no anonymous writes.)
  \item \textbf{T}ampering: can a caller modify data in transit or at rest?
        (Control: TLS everywhere; integrity checks on stored payloads.)
  \item \textbf{R}epudiation: can an actor deny an action?
        (Control: immutable audit logs with actor identity.)
  \item \textbf{I}nformation disclosure: can a caller read data beyond its
        authorization? (Control: object-level authorization on every
        resource fetch; minimal response fields.)
  \item \textbf{D}enial of service: can a caller exhaust capacity?
        (Control: rate limits, quotas, payload-size caps.)
  \item \textbf{E}levation of privilege: can a low-privilege caller perform
        high-privilege operations? (Control: server-side role checks; never
        trust client-supplied roles.)
\end{itemize}

Chapter~\ref{ch:api_security} turns each row into executable tests.

\paragraph{Ethics and responsible use.}
APIs move power: they decide whose data flows where, at what price, under
whose surveillance. The professional baseline adopted throughout this book:
collect the minimum data the operation requires; never log credentials,
tokens, or personal data; honor the provider's terms and the user's consent
when consuming; publish honest deprecation timelines when producing; and
treat synthetic-data discipline (never real PII in examples, tests, or
datasets) as non-negotiable. All datasets in this curriculum are fictional,
set in the Northwind Robotics universe.

\subsection{Scope Boundaries, Assumptions, and Risks}
\label{sec:ch0-scope}

\paragraph{Scope boundaries.}
This book teaches API engineering, not: frontend frameworks, database
administration, Kubernetes cluster operation (Chapter~\ref{ch:deployment}
uses it, does not teach it), language-agnostic theory for its own sake, or
vendor-specific cloud certifications. Where a vendor service is the
industry-standard reference (e.g.\ a managed gateway), we name it but
implement with local, offline equivalents.

\paragraph{Assumption register.}
Assumptions the rest of the book makes, stated once so later chapters can
build on them:

\begin{enumerate}[label=\textbf{A\arabic*.}]
  \item Python 3.11+ on Windows with PowerShell; commands are
        PowerShell-native.
  \item All network dependencies are simulated locally; no cloud accounts or
        API keys are available or required.
  \item Readers control a machine where they may open localhost ports and
        install packages into a virtual environment.
  \item Synthetic data only; the Northwind Robotics datasets in
        \texttt{datasets\textbackslash} are the canonical fixtures.
  \item Latency claims are always percentile claims; no claim is made from
        averages alone.
\end{enumerate}

\paragraph{Risk register (top entries).}

\begin{table}[h]
  \centering
  \small
  \begin{tabularx}{\textwidth}{X l l X}
    \toprule
    \textbf{Risk} & \textbf{Likelihood} & \textbf{Impact} & \textbf{Mitigation} \\
    \midrule
    Reader environment diverges from baseline & Medium & High &
      Sanity-check script (Listing~\ref{lst:ch0-sanity}) gates every lab. \\
    Percentile measurements misread due to tiny samples & High & Medium &
      Benchmark protocol fixes $n \geq 200$ and warmup discipline. \\
    Learners copy toy auth patterns into real systems & Medium & High &
      Warning boxes flag every non-production pattern in place. \\
    Offline stubs create false confidence about real services & Low & Medium &
      Chapters state exactly which behaviors the stub does not model. \\
    \bottomrule
  \end{tabularx}
  \caption{Initial risk register for the curriculum itself.}
  \label{tab:ch0-risk}
\end{table}

\paragraph{Success criteria template.}
Every chapter closes with measurable criteria of this form: \emph{Given
(prerequisites), when (procedure), then (observable outcome with a number).}
Example for this chapter: ``Given a fresh venv, when the quick-start lab is
executed per Section~\ref{sec:ch0-lab}, then the client completes 200
measured requests with zero errors and reports p50, p95, and p99, and the
learner can state which latency stage dominates on localhost.''

\paragraph{Terminology normalization.}
To keep twenty-one chapters unambiguous: \emph{endpoint} means one operation
on one path; \emph{resource} means a domain entity addressable by URI;
\emph{representation} means a serialized form of a resource;
\emph{consumer/provider} are roles, not machines (one service is often both);
\emph{error} is any non-success outcome; \emph{fault} is a defect that causes
errors; \emph{failure} is user-visible loss of service. Latency is always
client-observed unless marked \emph{server-side}.
%% ============================================================================
\section{Production-Ready Code Implementation}
\label{sec:ch0-code}

This section installs the two disciplines every later chapter assumes:
a \emph{reproducible environment} and a \emph{reproducible measurement
protocol}. All three listings are complete, typed, stdlib-only, and offline.

\subsection{Environment Setup Baseline}
\label{sec:ch0-env}

\paragraph{Minimal local stack (used from Chapter 1).}
Python 3.11+, a virtual environment, and the packages in this chapter's
appendix. From the repository root in PowerShell:

\begin{minted}{text}
python -m venv .venv
.\.venv\Scripts\Activate.ps1
python -m pip install --upgrade pip
pip install -r requirements.txt
\end{minted}

If activation is blocked by policy, run
\texttt{Set-ExecutionPolicy -Scope CurrentUser RemoteSigned} once, or invoke
\texttt{.\textbackslash .venv\textbackslash Scripts\textbackslash python.exe}
directly without activating.

\paragraph{Cloud-ready stack (used from Chapter 17).}
The local stack plus Docker Desktop (for container labs), a local Kubernetes
distribution, and a tunnel-free mindset: all ``cloud'' deployments in this
book target local clusters and local registries so the exercises remain
offline. The cloud-ready stack adds no Python dependencies beyond the
minimal stack.

\paragraph{Sanity check.}
Run Listing~\ref{lst:ch0-sanity} once per machine and once after any
environment change. It verifies the interpreter version, importability of
required packages, and availability of the localhost ports used by the labs
--- and it never touches the network beyond loopback.

\begin{minted}{python}
"""Environment sanity check for the API curriculum (offline, stdlib-only).

Verifies: interpreter version, required packages, localhost port
availability, and writes a structured report via logging.
Run:  python sanity_check.py
Exit code 0 means the environment is ready for every lab in the book.
"""
from __future__ import annotations

import importlib.util
import logging
import socket
import sys
from dataclasses import dataclass, field

MIN_PYTHON: tuple[int, int] = (3, 11)
REQUIRED_PACKAGES: tuple[str, ...] = (
    "fastapi", "uvicorn", "httpx", "pydantic", "pytest",
)
REQUIRED_PORTS: tuple[int, ...] = (8000, 8001, 8765)

logging.basicConfig(
    level=logging.INFO,
    format="%(levelname)-7s %(message)s",
)
logger = logging.getLogger("sanity_check")


@dataclass
class CheckReport:
    """Accumulates named check results for a final verdict."""

    failures: list[str] = field(default_factory=list)

    def record(self, name: str, ok: bool, detail: str) -> None:
        status = "PASS" if ok else "FAIL"
        logger.info("[%s] %s: %s", status, name, detail)
        if not ok:
            self.failures.append(name)


def check_python_version(report: CheckReport) -> None:
    current = sys.version_info[:2]
    ok = current >= MIN_PYTHON
    detail = f"found {current[0]}.{current[1]}, require >= {MIN_PYTHON[0]}.{MIN_PYTHON[1]}"
    report.record("python-version", ok, detail)


def check_packages(report: CheckReport) -> None:
    for package in REQUIRED_PACKAGES:
        found = importlib.util.find_spec(package) is not None
        detail = "importable" if found else "NOT INSTALLED"
        report.record(f"package:{package}", found, detail)


def check_port_available(port: int) -> tuple[bool, str]:
    """A port is 'available' if we can bind it on loopback right now."""
    try:
        with socket.socket(socket.AF_INET, socket.SOCK_STREAM) as probe:
            probe.setsockopt(socket.SOL_SOCKET, socket.SO_REUSEADDR, 1)
            probe.bind(("127.0.0.1", port))
        return True, "free"
    except OSError as exc:
        return False, f"unavailable ({exc.strerror or exc})"


def check_ports(report: CheckReport) -> None:
    for port in REQUIRED_PORTS:
        ok, detail = check_port_available(port)
        report.record(f"port:{port}", ok, detail)


def main() -> int:
    report = CheckReport()
    check_python_version(report)
    check_packages(report)
    check_ports(report)
    if report.failures:
        logger.error("Environment NOT ready; failed checks: %s",
                     ", ".join(report.failures))
        return 1
    logger.info("Environment ready for all labs.")
    return 0


if __name__ == "__main__":
    raise SystemExit(main())
\end{minted}
\label{lst:ch0-sanity}

\subsection{Evaluation Primer: Measuring Latency Correctly}
\label{sec:ch0-eval}

\paragraph{Why averages lie.}
Latency distributions of network services are heavy-tailed: most requests are
fast, a few are catastrophically slow, and the tail is where user pain lives.
The arithmetic mean of a heavy-tailed distribution is dominated by rare
outliers in small samples and hides them in large ones --- either way it
describes no actual user's experience. The standard remedy is
\emph{percentiles}: the value below which a given fraction of observations
falls. We always report three:

\begin{definition}[Latency percentile]
\label{def:percentile}
Given $n$ sorted latency samples $t_{(1)} \le t_{(2)} \le \dots \le t_{(n)}$,
the $q$-th percentile $p_q$ is $t_{(\lceil q n / 100 \rceil)}$. Thus $p_{50}$
(the median) is the typical request, $p_{95}$ is what the slowest 5\%
experience at best, and $p_{99}$ exposes the tail that averages conceal.
\end{definition}

\paragraph{Throughput vs.\ latency.}
Throughput $\lambda$ (requests/second) and latency $L$ are not independent.
By Little's Law, concurrency $C = \lambda \cdot L$; as load approaches
capacity, queueing theory predicts latency grows non-linearly even while
throughput still rises. Benchmark claims must therefore state the offered
load: ``p99 = 45\,ms at 200\,rps'' is a measurement; ``p99 = 45\,ms'' is
marketing. Throughout this book, benchmarks pin the load first.

\paragraph{Minimal benchmark protocol.}
Every performance claim in this book is produced by this protocol:

\begin{enumerate}
  \item \textbf{Warmup:} issue 50 unmeasured requests to stabilize caches,
        JIT-like effects, and connection pools.
  \item \textbf{Measure:} issue $n \geq 200$ sequential requests, timing each
        with \texttt{time.perf\_counter\_ns()} on the client.
  \item \textbf{Record:} keep raw samples; never aggregate before storing.
  \item \textbf{Report:} p50, p95, p99, max, error count, achieved
        throughput, and the exact configuration (variant V1--V4, payload
        size, machine class).
  \item \textbf{Repeat:} results quoted in prose are medians of 3 runs.
\end{enumerate}

\begin{warningbox}
Two measurement sins invalidate benchmarks silently. First, measuring with
\texttt{time.time()} or \texttt{datetime.now()}: wall clocks can jump
backwards under NTP corrections; always use a monotonic clock
(\texttt{time.perf\_counter\_ns}). Second, timing server-side only: the user
experiences the client-observed latency, which includes queueing and network
--- server logs systematically under-report it.
\end{warningbox}

\paragraph{The benchmark harness.}
Listing~\ref{lst:ch0-bench} is the harness used (and extended) throughout the
book. It is transport-agnostic: pass any callable and it returns a typed
\texttt{BenchmarkResult}. It deliberately uses only the standard library so
it runs in every chapter's environment.

\begin{minted}{python}
"""Reusable latency benchmark harness (stdlib-only, deterministic protocol).

Implements the book's minimal benchmark protocol: warmup, timed samples via
a monotonic clock, raw-sample retention, percentile reporting. No network
code lives here; the caller supplies a `send` callable so the same harness
measures urllib clients, httpx clients, or in-process fakes.

Run the built-in self-test:  python bench.py --selftest
"""
from __future__ import annotations

import argparse
import logging
import math
import statistics
import time
from collections.abc import Callable
from dataclasses import dataclass, field

logger = logging.getLogger("bench")

MIN_SAMPLES = 200
DEFAULT_WARMUP = 50


class BenchmarkError(RuntimeError):
    """Raised when the benchmark configuration is invalid."""


@dataclass(frozen=True)
class BenchmarkResult:
    """Immutable outcome of one benchmark run."""

    samples_ms: tuple[float, ...]
    errors: int
    wall_seconds: float
    label: str = ""

    @property
    def n(self) -> int:
        return len(self.samples_ms)

    @property
    def throughput_rps(self) -> float:
        if self.wall_seconds <= 0.0:
            raise BenchmarkError("wall time must be positive")
        return (self.n + self.errors) / self.wall_seconds

    def percentile(self, q: float) -> float:
        """Nearest-rank percentile per the book's Definition."""
        if not self.samples_ms:
            raise BenchmarkError("no samples collected")
        if not 0.0 < q <= 100.0:
            raise BenchmarkError(f"q must be in (0, 100], got {q}")
        ordered = sorted(self.samples_ms)
        rank = max(1, math.ceil(q * len(ordered) / 100.0))
        return ordered[rank - 1]

    def summary(self) -> dict[str, float | int | str]:
        return {
            "label": self.label,
            "n": self.n,
            "errors": self.errors,
            "p50_ms": round(self.percentile(50), 3),
            "p95_ms": round(self.percentile(95), 3),
            "p99_ms": round(self.percentile(99), 3),
            "max_ms": round(max(self.samples_ms), 3),
            "throughput_rps": round(self.throughput_rps, 1),
        }


def run_benchmark(
    send: Callable[[], None],
    *,
    samples: int = MIN_SAMPLES,
    warmup: int = DEFAULT_WARMUP,
    label: str = "",
) -> BenchmarkResult:
    """Execute the warmup + measured-runs protocol against `send`.

    `send` must perform exactly one request and raise on failure; raised
    exceptions are counted as errors and excluded from latency samples.
    """
    if samples < MIN_SAMPLES:
        raise BenchmarkError(
            f"samples must be >= {MIN_SAMPLES} for stable percentiles")
    if warmup < 0:
        raise BenchmarkError("warmup must be >= 0")

    for _ in range(warmup):
        try:
            send()
        except Exception:  # noqa: BLE001 -- warmup failures are tolerated
            logger.debug("warmup request failed", exc_info=True)

    timings: list[float] = []
    errors = 0
    start_ns = time.perf_counter_ns()
    for _ in range(samples):
        t0 = time.perf_counter_ns()
        try:
            send()
        except Exception:  # noqa: BLE001 -- errors are data, not crashes
            errors += 1
            logger.debug("measured request failed", exc_info=True)
            continue
        timings.append((time.perf_counter_ns() - t0) / 1_000_000.0)
    wall_seconds = (time.perf_counter_ns() - start_ns) / 1_000_000_000.0

    return BenchmarkResult(
        samples_ms=tuple(timings), errors=errors,
        wall_seconds=wall_seconds, label=label,
    )


def _selftest() -> int:
    """Deterministic offline self-test against an in-process fake."""
    state = {"calls": 0}

    def fake_send() -> None:
        state["calls"] += 1
        time.sleep(0.0005)  # deterministic-ish work, no wall-clock asserts
        if state["calls"] % 97 == 0:
            raise ConnectionError("simulated transient failure")

    result = run_benchmark(fake_send, samples=200, warmup=10,
                           label="selftest")
    report = result.summary()
    for key, value in report.items():
        logger.info("%-16s %s", key, value)
    assert result.n + result.errors == 200, "sample accounting broken"
    assert report["p50_ms"] <= report["p95_ms"] <= report["p99_ms"]
    assert result.errors > 0, "expected simulated errors"
    logger.info("selftest OK")
    return 0


if __name__ == "__main__":
    logging.basicConfig(level=logging.INFO, format="%(levelname)-7s %(message)s")
    parser = argparse.ArgumentParser(description=__doc__)
    parser.add_argument("--selftest", action="store_true")
    args = parser.parse_args()
    if args.selftest:
        raise SystemExit(_selftest())
    parser.error("pass --selftest (this module is a library otherwise)")
\end{minted}
\label{lst:ch0-bench}

Note the deliberate design choices: errors are \emph{counted, not raised},
because a benchmark that crashes on the first failure measures nothing;
exceptions are caught broadly only at this boundary, logged with
\texttt{exc\_info}, and converted into data; and the self-test asserts
ordering invariants ($p_{50} \le p_{95} \le p_{99}$) rather than absolute
timings, keeping it deterministic and machine-independent.

\subsection{A Complete Offline API and Client in the Standard Library}
\label{sec:ch0-miniapi}

Before frameworks, you must see the primitives. Listing~\ref{lst:ch0-miniapi}
implements a working JSON API with \texttt{http.server} --- the Northwind
Robotics parts catalog --- plus a client built on \texttt{urllib}. It
enforces the disciplines this book demands everywhere: explicit timeouts on
every call, structured JSON errors, input validation with specific
exceptions, and logging instead of \texttt{print}. Chapter~\ref{ch:fastapi}
will replace the plumbing with FastAPI; the disciplines remain.

\begin{minted}{python}
"""Northwind Robotics parts API + client, stdlib-only and fully offline.

Server:  GET /healthz                 -> {"status": "ok"}
         GET /v1/parts                -> {"items": [...]}
         GET /v1/parts/<sku>          -> part object or RFC-shaped error
Client:  typed functions with explicit timeouts and error mapping.

Run server:  python mini_api.py serve --port 8000
Run client demo (spawns server in-process):  python mini_api.py demo
"""
from __future__ import annotations

import argparse
import json
import logging
import threading
import urllib.error
import urllib.request
from dataclasses import dataclass
from http.server import BaseHTTPRequestHandler, ThreadingHTTPServer
from typing import Any

logger = logging.getLogger("mini_api")

DEFAULT_HOST = "127.0.0.1"
DEFAULT_PORT = 8000
CLIENT_TIMEOUT_S = 2.0  # every network call in this book carries a timeout

# Synthetic fixture data (Northwind Robotics universe). No real PII.
PARTS: dict[str, dict[str, Any]] = {
    "NR-AXLE-001": {"sku": "NR-AXLE-001", "name": "Drive axle, Mk IV",
                    "unit_price_cents": 4599, "stock": 37},
    "NR-LIDAR-204": {"sku": "NR-LIDAR-204", "name": "LiDAR puck, 16-channel",
                     "unit_price_cents": 189900, "stock": 4},
    "NR-GRIP-011": {"sku": "NR-GRIP-011", "name": "Parallel gripper jaw",
                    "unit_price_cents": 12350, "stock": 120},
}


def problem(status: int, title: str, detail: str) -> bytes:
    """RFC 9457-flavored error body (formalized in Chapter 7)."""
    body = {
        "type": "about:blank",
        "title": title,
        "status": status,
        "detail": detail,
    }
    return json.dumps(body).encode("utf-8")


class PartsHandler(BaseHTTPRequestHandler):
    server_version = "NorthwindParts/0.1"

    def log_message(self, fmt: str, *args: Any) -> None:
        logger.info("%s - %s", self.address_string(), fmt % args)

    def _send_json(self, status: int, payload: bytes) -> None:
        self.send_response(status)
        self.send_header("Content-Type", "application/json")
        self.send_header("Content-Length", str(len(payload)))
        self.end_headers()
        self.wfile.write(payload)

    def do_GET(self) -> None:  # noqa: N802 -- stdlib naming contract
        if self.path == "/healthz":
            self._send_json(200, json.dumps({"status": "ok"}).encode())
            return
        if self.path == "/v1/parts":
            payload = json.dumps({"items": list(PARTS.values())}).encode()
            self._send_json(200, payload)
            return
        if self.path.startswith("/v1/parts/"):
            sku = self.path.removeprefix("/v1/parts/").strip()
            part = PARTS.get(sku)
            if part is None:
                self._send_json(404, problem(
                    404, "Not Found", f"no part with sku '{sku}'"))
                return
            self._send_json(200, json.dumps(part).encode())
            return
        self._send_json(404, problem(
            404, "Not Found", f"unknown path '{self.path}'"))


def make_server(host: str = DEFAULT_HOST,
                port: int = DEFAULT_PORT) -> ThreadingHTTPServer:
    return ThreadingHTTPServer((host, port), PartsHandler)


# ------------------------------- client side -------------------------------

class ApiClientError(RuntimeError):
    """Base class for client-side failures."""


class ApiHttpError(ApiClientError):
    """The server answered with a non-2xx status."""

    def __init__(self, status: int, detail: str) -> None:
        super().__init__(f"HTTP {status}: {detail}")
        self.status = status
        self.detail = detail


@dataclass(frozen=True)
class PartsClient:
    """Minimal typed client for the parts API."""

    base_url: str = f"http://{DEFAULT_HOST}:{DEFAULT_PORT}"
    timeout_s: float = CLIENT_TIMEOUT_S

    def _get_json(self, path: str) -> Any:
        if not path.startswith("/"):
            raise ValueError(f"path must start with '/', got {path!r}")
        request = urllib.request.Request(self.base_url + path, method="GET")
        try:
            with urllib.request.urlopen(request,
                                        timeout=self.timeout_s) as resp:
                return json.loads(resp.read().decode("utf-8"))
        except urllib.error.HTTPError as exc:
            detail = exc.read().decode("utf-8", errors="replace")
            raise ApiHttpError(exc.code, detail) from exc
        except (urllib.error.URLError, TimeoutError, OSError) as exc:
            raise ApiClientError(f"transport failure: {exc}") from exc

    def health(self) -> bool:
        return self._get_json("/healthz")["status"] == "ok"

    def list_parts(self) -> list[dict[str, Any]]:
        return list(self._get_json("/v1/parts")["items"])

    def get_part(self, sku: str) -> dict[str, Any]:
        if not sku:
            raise ValueError("sku must be non-empty")
        return dict(self._get_json(f"/v1/parts/{sku}"))


def demo() -> int:
    """Start the server in-process and exercise the client. Offline."""
    server = make_server()
    thread = threading.Thread(target=server.serve_forever, daemon=True)
    thread.start()
    try:
        client = PartsClient()
        logger.info("health: %s", client.health())
        parts = client.list_parts()
        logger.info("catalog size: %d", len(parts))
        part = client.get_part("NR-LIDAR-204")
        logger.info("fetched: %s (%d cents)",
                    part["name"], part["unit_price_cents"])
        try:
            client.get_part("NO-SUCH-SKU")
        except ApiHttpError as exc:
            logger.info("expected 404 mapped cleanly: status=%d", exc.status)
    finally:
        server.shutdown()
        server.server_close()
        thread.join(timeout=2.0)
    return 0


def main() -> int:
    logging.basicConfig(level=logging.INFO,
                        format="%(levelname)-7s %(message)s")
    parser = argparse.ArgumentParser(description=__doc__)
    sub = parser.add_subparsers(dest="command", required=True)
    serve = sub.add_parser("serve", help="run the server in the foreground")
    serve.add_argument("--port", type=int, default=DEFAULT_PORT)
    sub.add_parser("demo", help="in-process server + client round trip")
    args = parser.parse_args()

    if args.command == "serve":
        if not 1024 <= args.port <= 65535:
            parser.error("port must be in [1024, 65535]")
        server = make_server(port=args.port)
        logger.info("serving on http://%s:%d", DEFAULT_HOST, args.port)
        try:
            server.serve_forever()
        except KeyboardInterrupt:
            logger.info("shutting down")
        finally:
            server.server_close()
        return 0
    return demo()


if __name__ == "__main__":
    raise SystemExit(main())
\end{minted}
\label{lst:ch0-miniapi}

\begin{keyconceptbox}
Every listing in this book obeys five rules, visible already in these three:
(1)~explicit timeouts on every network call; (2)~typed function signatures
and dataclasses; (3)~specific exceptions with chained causes
(\texttt{raise ... from exc}); (4)~logging instead of \texttt{print};
(5)~offline determinism --- no listing requires the internet, a cloud
account, or an API key.
\end{keyconceptbox}
%% ============================================================================
\section{Common Pitfalls \& Anti-Patterns}
\label{sec:ch0-pitfalls}

These are the failure patterns this book's authors see most often in code
reviews and incident post-mortems. Each is stated with its detection signal
and its fix; the chapters referenced contain the full treatment.

\subsection{Consumer-Side Pitfalls}

\begin{description}
  \item[Missing timeouts (the cardinal sin).] Any network call without a
    timeout can block \emph{forever}; TCP's default retransmission behavior
    happily waits minutes. Detection: thread dumps full of blocked I/O
    threads; fix: a connect \emph{and} read timeout on every call
    (Section~\ref{sec:ch0-warstory} shows what happens otherwise).
  \item[Retrying everything, immediately, forever.] Unbounded retries turn a
    partial outage into a self-inflicted denial of service. Fix: bounded
    retries with exponential backoff and jitter, only on retryable
    (idempotent or explicitly safe) operations (Chapters
    \ref{ch:consuming_python}, \ref{ch:idempotency}, \ref{ch:rate_limiting}).
  \item[Parsing errors by string-matching.] Matching on
    \texttt{"something went wrong"} in a response body is a contract bug
    waiting to be renamed. Fix: branch on status codes and
    machine-readable error types (Chapter~\ref{ch:problem_details}).
  \item[Ignoring the contract until runtime.] Discovering a renamed field in
    production is a process failure. Fix: schema validation in CI
    (Chapters \ref{ch:serialization}, \ref{ch:testing}).
  \item[Secrets in source or logs.] API keys committed to git or printed in
    debug logs are the most common real-world API breach vector. Fix:
    environment variables, secret stores, log redaction
    (Chapter~\ref{ch:client_security}).
\end{description}

\subsection{Provider-Side Pitfalls}

\begin{description}
  \item[Leaky abstraction responses.] Returning stack traces, ORM errors, or
    internal hostnames in error bodies leaks your architecture to attackers.
    Fix: RFC~9457 problem details with safe, curated fields.
  \item[Breaking changes in place.] ``It's an internal API'' is how
    distributed monoliths are born; consumers couple silently. Fix: version
    every contract and automate compatibility checks
    (Chapter~\ref{ch:versioning}).
  \item[Offset pagination on mutating collections.] Items inserted while a
    consumer pages shift every window, duplicating or skipping records. Fix:
    cursor-based pagination (Chapter~\ref{ch:pagination}).
  \item[200 OK with an error inside.] Wrapping failures in success envelopes
    defeats HTTP intermediaries, caches, and every monitoring tool. Fix:
    correct status codes; bodies carry detail, statuses carry semantics.
  \item[No resource limits.] Unbounded page sizes, payload sizes, and query
    complexity are invitations to accidental (or deliberate) denial of
    service. Fix: caps plus rate limiting (Chapter~\ref{ch:rate_limiting}).
\end{description}

\subsection{Organizational Anti-Patterns}

\begin{description}
  \item[The API as afterthought.] Designing the UI first and exposing
    whatever it needed produces chatty, anemic APIs. Fix: contract-first
    design reviews.
  \item[Docs as a wiki page.] Unversioned prose documentation decays within
    weeks. Fix: machine-readable contracts as the source of truth, docs
    generated from them \cite{openapi31}.
  \item[No owner.] An API owned by ``the team'' is owned by no one;
    deprecation, incidents, and consumer questions all stall. Fix: a named
    owner per contract (governance baseline,
    Section~\ref{sec:ch0-governance}).
\end{description}

\subsection{Common Myths About APIs}
\label{sec:ch0-myths}

\begin{description}
  \item[``If it returns JSON over HTTP, it is REST.''] False. REST is a set
    of architectural constraints \cite{fielding2000rest}; most
    JSON-over-HTTP APIs are RPC-flavored HTTP interfaces. That is often fine
    --- but call things by their names, because the constraints you actually
    honor determine which guarantees (cacheability, statelessness) you may
    assume.
  \item[``HTTP 200 means it worked.''] It means the server produced a
    response it considers successful. Semantic failures ride inside 200
    bodies all the time; monitoring that counts only status codes is blind
    to them.
  \item[``Internal APIs don't need contracts or versioning.''] Internal
    consumers are still consumers, and they couple silently. The blast
    radius of an unversioned internal API is often \emph{larger} than a
    public one because nobody tracks who depends on it.
  \item[``Timeouts make systems slower.''] Timeouts cap the blast radius of
    slow dependencies; the slow path exists whether or not you bound it.
    The unbounded version is strictly worse (Section~\ref{sec:ch0-warstory}).
  \item[``Retries fix flakiness.''] Retries fix \emph{transient transport}
    failures and amplify everything else: they multiply load during
    outages and duplicate non-idempotent effects. A retry policy is a
    designed artifact, not a default.
  \item[``The gateway will handle security.''\/] Gateways terminate TLS and
    enforce coarse policies; they cannot fix broken object-level
    authorization or injection inside business logic. Security is layered,
    and the innermost layer is your code (Chapter~\ref{ch:api_security}).
  \item[``Average latency is a fine health metric.''] Averages hide the
    tail where user pain lives. Percentiles under stated load, or it did
    not happen (Section~\ref{sec:ch0-eval}).
  \item[``Documentation keeps itself up to date.''] Only generated
    artifacts stay honest. Hand-written prose decays; machine-readable
    contracts with docs rendered from them do not \cite{openapi31}.
\end{description}

\subsection{Common Mistakes \& Debugging Gallery}
\label{sec:ch0-gallery}

A field guide to symptoms you will meet in the labs:

\begin{table}[h]
  \centering
  \small
  \begin{tabularx}{\textwidth}{X X X}
    \toprule
    \textbf{Symptom} & \textbf{Likely cause} & \textbf{First check} \\
    \midrule
    \texttt{ConnectionRefusedError} on localhost &
      Server not running, or bound to a different port/interface &
      \texttt{Test-NetConnection 127.0.0.1 -Port 8000} in PowerShell. \\
    Request hangs, no exception &
      Missing client timeout; server accepted but never answered &
      Add/verify connect + read timeouts; check server logs. \\
    Intermittent \texttt{RemoteDisconnected} &
      Server closed an idle keep-alive connection the client reused &
      Retry once on idempotent methods; shorten client keep-alive. \\
    HTTP 415 Unsupported Media Type &
      Missing or wrong \texttt{Content-Type} header &
      Compare request headers against the contract's media type. \\
    JSON decode error on error paths &
      Error body is not JSON (proxy HTML page, empty body) &
      Log status + first 200 bytes of body before decoding. \\
    p99 spike with flat p50 &
      Queueing at saturation, GC pauses, or one slow downstream &
      Decompose per Equation~\eqref{eq:ch0-budget}; check pool metrics. \\
    Works in demo, fails under load test &
      Connection exhaustion, thread-pool starvation, retry storms &
      Count open sockets; check retry policy and pool sizing. \\
    \bottomrule
  \end{tabularx}
  \caption{Symptom-to-cause gallery for the most common lab failures.}
  \label{tab:ch0-gallery}
\end{table}

%% ============================================================================
\section{Hands-On Production Lab \& Real-World Case Study}
\label{sec:ch0-lab}

\subsection{Quick-Start Lab: Build, Run, Measure (2--3 Hours, Offline)}

\textbf{Objective.} Stand up the reference stack at variant V1+V2, run the
benchmark protocol against it, and produce your first honest percentile
report.

\textbf{Prerequisites.} Sanity check green
(Listing~\ref{lst:ch0-sanity}).

\paragraph{Step 1 --- Environment (15 min).}
Create and activate the virtual environment and install dependencies exactly
as in Section~\ref{sec:ch0-env}. Run the sanity check and keep its output;
labs in later chapters assume it.

\paragraph{Step 2 --- Save the code (15 min).}
Save Listing~\ref{lst:ch0-miniapi} as \texttt{mini\_api.py} and
Listing~\ref{lst:ch0-bench} as \texttt{bench.py} in a lab folder. Verify the
harness self-test:

\begin{minted}{text}
python bench.py --selftest
\end{minted}

Expect the summary keys (\texttt{p50\_ms}, \texttt{p95\_ms},
\texttt{p99\_ms}, \texttt{errors}, \texttt{throughput\_rps}) and a final
\texttt{selftest OK}.

\paragraph{Step 3 --- Run the server (10 min).}
In one PowerShell window:

\begin{minted}{text}
python mini_api.py serve --port 8000
\end{minted}

In a second window, smoke-test with the stdlib client:

\begin{minted}{text}
python mini_api.py demo
\end{minted}

\paragraph{Step 4 --- Measure (45 min).}
Write \texttt{lab0\_measure.py} that imports both modules and measures the
real client against the real server:

\begin{minted}{python}
"""Lab 0 measurement driver: benchmark the stdlib client end to end."""
from __future__ import annotations

import logging
import threading

from bench import run_benchmark
from mini_api import PartsClient, make_server

logger = logging.getLogger("lab0")


def main() -> int:
    logging.basicConfig(level=logging.INFO,
                        format="%(levelname)-7s %(message)s")
    server = make_server(port=8000)
    thread = threading.Thread(target=server.serve_forever, daemon=True)
    thread.start()
    try:
        client = PartsClient(base_url="http://127.0.0.1:8000")
        result = run_benchmark(
            lambda: client.get_part("NR-AXLE-001"),
            samples=500, warmup=50, label="stdlib-local",
        )
        for key, value in result.summary().items():
            logger.info("%-16s %s", key, value)
        if result.errors != 0:
            logger.error("expected zero errors on localhost, got %d",
                         result.errors)
            return 1
    finally:
        server.shutdown()
        server.server_close()
        thread.join(timeout=2.0)
    return 0


if __name__ == "__main__":
    raise SystemExit(main())
\end{minted}
\label{lst:ch0-labdriver}

\paragraph{Step 5 --- Analyze (45 min).}
Answer in writing: (a)~Which latency stage of
Equation~\eqref{eq:ch0-budget} dominates on localhost, and why does
$L_{\mathrm{TLS}}$ not appear? (b)~What is the ratio $p_{99}/p_{50}$? What
does a ratio above 3 suggest? (c)~Repeat with the server under a second
concurrent client; how do p50 and p99 move relative to each other?

\paragraph{Step 6 --- Break it deliberately (30 min).}
Introduce one controlled fault at a time and record which taxonomy category
(Table~\ref{tab:ch0-taxonomy}) each produces: stop the server mid-benchmark
(transport); request a nonexistent SKU and count 404s (contract/semantic
boundary --- argue which); remove the client timeout and kill the server
before answering (transport, and observe the hang); send a malformed path
(contract). This fault-injection habit is the seed of Chapter 19's chaos
testing.

\paragraph{Acceptance criteria.}
\begin{itemize}
  \item[$\square$] Sanity check reports all PASS.
  \item[$\square$] 500-request benchmark completes with zero errors.
  \item[$\square$] Written answers name the dominant latency stage and the
    $p_{99}/p_{50}$ ratio.
  \item[$\square$] Fault-injection table maps each induced failure to a
    taxonomy category.
\end{itemize}

\subsection{Production War Story: The Missing Timeout}
\label{sec:ch0-warstory}

Names are fictionalized; the failure pattern is one of the most common in the
industry \cite{nygard2018releaseit}.

\paragraph{Background.}
Northwind Robotics' fleet dashboard (consumer) called the inventory service
(provider) over HTTP to render each robot's parts manifest. The client
library defaulted to \emph{no timeout} --- ``we'll add timeouts later'' had
been on the backlog for two quarters. The inventory service, in turn,
depended on a catalog database.

\paragraph{Trigger.}
At 14:07 on a Tuesday, a routine database failover made catalog queries slow:
not failing, just slow --- 30 seconds instead of 30 milliseconds. Slow
failures are the dangerous kind; nothing alarms because nothing has
``failed.''

\paragraph{Cascade.}
Inventory-service workers blocked on catalog queries. Its thread pool (64
threads) exhausted in under a minute, so inventory requests from the
dashboard queued, then blocked. Dashboard server workers, holding open
requests with no timeout, blocked in turn; its pool (128 threads) exhausted
four minutes later. The load balancer, seeing all dashboard workers busy,
began returning 503s for \emph{every} route --- including
\texttt{/healthz} and the static marketing pages that had no dependency on
inventory at all. By 14:19, a slow SELECT statement had taken down the
entire customer-facing site. Retry logic in a third system (the mobile app's
aggressive immediate-retry policy) tripled the offered load during the
incident, deepening the queue.

\paragraph{Detection and recovery.}
The on-call engineer's percentile dashboards told the story in one glance:
p50 dashboard latency flat until 14:07, then p99 vertical. Recovery was
manual thread-pool drainage via rolling restarts; the catalog failover
completed on its own at 14:31. Total customer-visible outage: 24 minutes ---
more than half a 99.9\% monthly error budget
(Equation~\eqref{eq:ch0-errorbudget}) spent in one afternoon.

\paragraph{Root cause and fixes.}
The post-mortem classified the incident across three taxonomy categories: an
\emph{operational} trigger (slow dependency), amplified by a \emph{transport}
defect (missing timeouts --- no failure semantics defined between the
systems), worsened by a \emph{semantic} one (immediate unbounded retries).
Fixes, in order of leverage: (1)~connect and read timeouts on every call
(500\,ms / 2\,s); (2)~bounded retries with exponential backoff and jitter on
idempotent reads only; (3)~bulkheads --- separate thread pools per
dependency so one slow dependency cannot starve unrelated routes;
(4)~circuit breakers to fail fast while the dependency recovers; (5)~a load
test reproducing the scenario in CI. Chapters \ref{ch:consuming_python},
\ref{ch:idempotency}, \ref{ch:rate_limiting}, and \ref{ch:microservices}
implement each fix; Chapter~\ref{ch:observability} builds the dashboards that
caught it.

\begin{warningbox}
The default timeout of most HTTP client libraries is \emph{infinite} or
absurdly long. Treat any client configuration without an explicit timeout as
a Sev-2 defect. There is no safe default; there is only your latency budget.
\end{warningbox}

\subsection{Diagnostic Playbook: ``The API Is Slow'' / ``The API Is Broken''}
\label{sec:ch0-playbook}

When the page arrives, run this sequence before forming theories:

\begin{enumerate}
  \item \textbf{Define the symptom precisely.} Which operation, which
    consumers, since when? Get a failing request ID or a reproducible call.
    ``Slow'' and ``broken'' are different playbooks from here.
  \item \textbf{Reproduce from a neutral client.} A single measured call from
    outside the complaining system (your Chapter 2 client) isolates
    consumer-side causes.
  \item \textbf{Classify into the taxonomy.} Transport (cannot connect,
    TLS errors, 502s)? Contract (400/422, schema mismatches)? Semantic
    (200 with wrong data)? Operational (latency tail, saturation)? Security
    (401/403, blocked IPs)? Table~\ref{tab:ch0-taxonomy} maps each to its
    detection strategy.
  \item \textbf{Localize with the budget.} Decompose latency per
    Equation~\eqref{eq:ch0-budget}; the stage that exceeds its budget is
    your suspect. For errors, walk the request path hop by hop.
  \item \textbf{Check the tail, not the mean.} p99 first; if p99 is fine,
    the complaint is about a specific consumer segment, not the fleet.
  \item \textbf{Form one hypothesis, test one change.} Change a single
    variable, re-measure with the benchmark protocol, keep the raw samples.
  \item \textbf{Write it down.} Every incident ends with the taxonomy
    classification, the fix, and the test that prevents recurrence.
\end{enumerate}
%% ============================================================================
\section{Self-Assessment Exercises \& Deep-Dive Questions}
\label{sec:ch0-selfassess}

\subsection{Readiness Self-Assessment}

Score yourself honestly (0 = no idea, 1 = heard of it, 2 = can do it) before
starting Chapter 1. Total $\geq 14$ means you are ready; any 0 identifies a
topic to shore up first.

\begin{table}[h]
  \centering
  \small
  \begin{tabularx}{\textwidth}{X c}
    \toprule
    \textbf{Skill} & \textbf{Score (0--2)} \\
    \midrule
    Create and activate a Python virtual environment in PowerShell & \\
    Write a typed Python function with a dataclass return value & \\
    Explain what happens between \texttt{urlopen} and receiving bytes & \\
    Read an HTTP request/response pair in raw text form & \\
    Branch on HTTP status codes and content types & \\
    Explain the difference between TCP and HTTP & \\
    Compute a percentile from a sorted list by hand & \\
    Use git to clone, branch, commit, and diff & \\
    Explain what TLS authenticates and encrypts & \\
    State one reason an average is a bad latency summary & \\
    \bottomrule
  \end{tabularx}
  \caption{Readiness self-assessment (target: total $\geq 14$, no zeros).}
  \label{tab:ch0-readiness}
\end{table}

\subsection{Exercises}

\begin{enumerate}[label=\textbf{E0.\arabic*.}]
  \item \textbf{[junior] Classify the incidents.} For each: (a)~a proxy
    returns 502 during a deploy; (b)~a field changes type from string to
    number; (c)~prices are returned in the wrong currency with HTTP 200;
    (d)~p99 triples every day at batch-job time; (e)~a customer reads
    another customer's invoice by incrementing an ID. Map each to the
    taxonomy and name its detection strategy.
  \item \textbf{[junior] Percentiles by hand.} Given sorted latencies (ms):
    4, 5, 5, 6, 7, 9, 12, 18, 40, 220 --- compute $p_{50}$, $p_{90}$, the
    mean, and explain which single number you would put on a dashboard.
  \item \textbf{[junior] Extend the mini API.} Add
    \texttt{GET /v1/parts?in\_stock=true} filtering to
    Listing~\ref{lst:ch0-miniapi}, returning a 400 problem body for unknown
    query parameters. Keep it stdlib-only and typed.
  \item \textbf{[mid] Error budget arithmetic.} Northwind's fleet API has a
    99.95\% monthly SLO. Compute the downtime budget in minutes for a
    30-day month, then convert it into ``allowed 5xx responses per million
    requests.'' Which unit is easier to alert on, and why?
  \item \textbf{[mid] Harden the benchmark.} Modify
    Listing~\ref{lst:ch0-bench} to interleave three payload sizes in a
    fixed, seeded pattern and report per-size percentiles. Why must the
    pattern be seeded rather than random?
  \item \textbf{[mid] Maturity audit.} Take the last API you worked on (or
    Listing~\ref{lst:ch0-miniapi}) and score it L0--L4. List the three
    cheapest changes that would raise it one level.
  \item \textbf{[senior] Budget decomposition.} Instrument
    Listing~\ref{lst:ch0-miniapi} to timestamp request receipt, handler
    start, and response flush; compute a per-stage latency table over 500
    requests and reconcile it with Equation~\eqref{eq:ch0-budget}. Which
    stage's budget would you tighten first?
  \item \textbf{[senior] Timeout theorem.} Using the war story's numbers,
    estimate the maximum request rate the dashboard could have sustained
    with 128 workers and 30-second dependency stalls. Then show how a
    2-second timeout changes the answer (Little's Law, $C=\lambda L$).
  \item \textbf{[senior] STRIDE the mini API.} Produce a six-row STRIDE
    table for Listing~\ref{lst:ch0-miniapi} as written. Which threats does
    the code actually control today? Which does it ignore?
  \item \textbf{[staff] Design the governance gate.} Draft the CI checks
    (as a checklist, not code) that would keep every Northwind API at L2
    minimum: what fails a build, what warns, and who can override. Compare
    with the addendum's CI quality-gate template.
\end{enumerate}

\subsection{Deep-Dive Questions}

\begin{enumerate}[label=\textbf{D0.\arabic*.}]
  \item Definition~\ref{def:api} makes failure semantics part of the
    contract. Argue for and against treating a change in retryability
    guidance (``clients may now retry POST X'') as a breaking change.
  \item The maturity model is linear; real organizations are not. Describe a
    plausible system that is simultaneously L4 in governance and L0 in
    resilience, and the risk that combination creates.
  \item Little's Law ($C = \lambda L$) assumes a stable system. Which
    assumptions break during a retry storm, and what does the law then
    mislead you about?
  \item Percentile latency SLOs can be gamed by shedding load (returning
    fast errors). Design an SLI definition that resists this gaming.
  \item Every lab in this book is offline. What does the offline constraint
    \emph{prevent} you from learning, and how would you compensate before
    touching a real provider?
\end{enumerate}

%% ============================================================================
\section{Interview Q\&A Vault}
\label{sec:ch0-qa}

Thirty-two questions, ordered junior $\to$ staff. Answers are intentionally
concise; interviews reward precision over volume.

\begin{enumerate}[label=\textbf{Q\arabic*.}, leftmargin=3.2em]
  \item \textbf{[junior] What is an API?}\\
    A versioned contract between provider and consumers specifying
    operations, representations, and failure semantics over a transport
    protocol (Definition~\ref{def:api}). Not the implementation --- the
    agreement.
  \item \textbf{[junior] What is the difference between a library and an
    API?}\\
    A library executes in your process and address space; an API crosses a
    process or network boundary, so it has transport, versioning, and
    partial-failure concerns a library never faces.
  \item \textbf{[junior] Name the parts of an HTTP request.}\\
    Request line (method, target, version), headers, optional body. The
    response mirrors this: status line, headers, optional body
    \cite{rfc9110}.
  \item \textbf{[junior] What do the HTTP status-code classes mean?}\\
    1xx informational, 2xx success, 3xx redirection, 4xx client fault,
    5xx server fault. The class alone localizes responsibility before you
    read a single body byte.
  \item \textbf{[junior] GET vs.\ POST?}\\
    GET retrieves a representation and must be safe and idempotent; POST
    submits an operation the server defines, with no such guarantee.
    Caches and retries treat them differently because the semantics
    differ.
  \item \textbf{[junior] What is JSON and why did it win over XML for web
    APIs?}\\
    A text serialization of objects, arrays, numbers, strings, booleans,
    and null. It won on parsing simplicity, direct mapping to programming
    language data structures, and payload size --- at the cost of no
    native schema or comment support, which Chapter 4 addresses.
  \item \textbf{[junior] What is a timeout and why does every call need
    one?}\\
    A bound on how long you wait for a connection or a response. Without
    it, a silent peer blocks your worker forever; at scale that exhausts
    pools and cascades (Section~\ref{sec:ch0-warstory}).
  \item \textbf{[junior] What is an idempotent operation?}\\
    One where performing it $n$ times has the same server-side effect as
    performing it once. GET, PUT, and DELETE are specified idempotent;
    POST is not, which is exactly why idempotency keys exist
    (Chapter~\ref{ch:idempotency}).
  \item \textbf{[junior] What is a status 404 vs.\ 410?}\\
    404: not found, possibly transient. 410: gone permanently, a deliberate
    deprecation signal. Returning the right one lets consumers automate
    their response.
  \item \textbf{[junior] What does TLS give an API connection?}\\
    Server authentication, confidentiality, and integrity in transit ---
    and with mutual TLS, client authentication too. It does nothing for
    authorization or input validation.
  \item \textbf{[mid] What is REST, precisely?}\\
    An architectural style \cite{fielding2000rest}: client--server,
    stateless, cacheable, uniform interface (resources identified by URIs,
    manipulated through representations, self-descriptive messages,
    hypermedia), layered system, optional code-on-demand. A JSON-over-HTTP
    CRUD endpoint is not automatically REST.
  \item \textbf{[mid] Explain the difference between authentication and
    authorization.}\\
    Authentication establishes \emph{who} the caller is; authorization
    decides \emph{what} that caller may do. The most damaging modern API
    flaw --- broken object-level authorization --- is an authenticated
    caller reaching objects it should not, i.e.\ an authorization failure.
  \item \textbf{[mid] Why do averages mislead for latency? What do you
    report instead?}\\
    Latency is heavy-tailed; the mean describes no real user and hides the
    tail where pain concentrates. Report percentiles: p50 for the typical
    experience, p95/p99 for the tail, always with sample size and offered
    load.
  \item \textbf{[mid] What is an API contract test?}\\
    A test that verifies provider behavior against the published contract
    (schema, status codes, error shapes), often generated from an OpenAPI
    document, and consumer tests that pin the fields the consumer actually
    uses. They run in CI so a breaking change fails a build, not a customer.
  \item \textbf{[mid] Your consumer gets HTTP 429. What happened and what do
    you do?}\\
    The provider rate-limited you. Read \texttt{Retry-After}, back off with
    exponential delay and jitter, reduce concurrency, and if the limit is
    chronic, negotiate quota or cache more. Never retry immediately
    (Chapter~\ref{ch:rate_limiting}).
  \item \textbf{[mid] What is connection pooling and why does it matter?}\\
    Reusing TCP/TLS connections across requests. Handshakes cost round
    trips (budget stage $L_{\mathrm{TLS}}$); pooling amortizes them,
    often the single largest latency win for chatty clients
    (Chapter~\ref{ch:consuming_python}).
  \item \textbf{[mid] When should you paginate, and offset vs.\ cursor?}\\
    Whenever a collection can outgrow one response. Offset pagination is
    simple but unstable under concurrent inserts and gets slower deeper in;
    cursor (keyset) pagination is stable and index-friendly. Default to
    cursors for large or mutating collections (Chapter~\ref{ch:pagination}).
  \item \textbf{[mid] What belongs in an error response body?}\\
    RFC~9457 fields: \texttt{type} (machine-readable URI), \texttt{title},
    \texttt{status}, \texttt{detail}, \texttt{instance}, plus curated
    extensions. Never stack traces, SQL, internal hostnames, or secrets.
  \item \textbf{[mid] Semantic versioning for APIs: what breaks what?}\\
    Removing or renaming fields, tightening types, changing semantics:
    breaking, new major. Adding optional fields or new operations:
    backward compatible, minor. Deprecating with a sunset date: policy,
    not schema (Chapter~\ref{ch:versioning}).
  \item \textbf{[mid] What is the N+1 problem in API composition?}\\
    One list call followed by $n$ detail calls. It multiplies latency and
    load; fix with batch endpoints, embedded expansions, or GraphQL-style
    selection --- chosen consciously, since each trades server complexity
    differently.
  \item \textbf{[senior] Design retries correctly. What is the full
    checklist?}\\
    (1)~Only retry on retryable failures: transport errors, 429, 5xx ---
    never 4xx semantics; (2)~only idempotent operations, or add
    idempotency keys; (3)~bounded attempts (typically 3); (4)~exponential
    backoff with full jitter; (5)~a retry budget so retries stay a small
    fraction of traffic; (6)~respect \texttt{Retry-After}; (7)~the total
    retry deadline must fit the caller's own latency budget.
  \item \textbf{[senior] Explain Little's Law and one operational use.}\\
    $C = \lambda L$: concurrency equals arrival rate times latency. Use:
    if a dependency slows from 30\,ms to 30\,s at 100\,rps, in-flight
    requests grow from 3 to 3000 --- you can compute the exact pool size
    at which you die, which is why timeouts and bulkheads are sized from
    this law.
  \item \textbf{[senior] What is a circuit breaker and what state machine
    does it run?}\\
    Closed (normal) $\to$ Open (failing fast) after an error threshold
    $\to$ Half-Open (probe with limited traffic) $\to$ Closed on success
    or Open on failure. It converts slow failures into fast ones,
    protecting both caller pools and the recovering dependency
    \cite{nygard2018releaseit}.
  \item \textbf{[senior] Error budget: what is it and how does it change
    behavior?}\\
    The allowed unreliability: $1 - \mathrm{SLO}$. At 99.9\% monthly, that
    is 43.2 minutes (Equation~\eqref{eq:ch0-errorbudget}). Budget
    remaining $\Rightarrow$ ship features; budget burning fast
    $\Rightarrow$ freeze releases, invest in reliability. It turns the
    reliability-vs-velocity argument into arithmetic \cite{beyer2016sre}.
  \item \textbf{[senior] How do you make a POST safely retryable?}\\
    Client generates an idempotency key per logical operation; server
    stores (key, response) with a TTL and replays the stored response on
    duplicate keys; unique constraints enforce exactly-once effects in
    storage. Watch the races: concurrent first-delivery and retry must
    serialize on the key (Chapter~\ref{ch:idempotency}).
  \item \textbf{[senior] GraphQL vs.\ REST vs.\ gRPC: how do you choose?}\\
    REST: public contracts, caching, broad interoperability. GraphQL:
    many clients with divergent data needs, aggregating several backends,
    at the cost of query-complexity governance. gRPC: internal
    service-to-service, strict schemas, streaming, and low latency, at
    the cost of browser and tooling friction. Choose per consumer, not
    per fashion (Chapters \ref{ch:graphql}--\ref{ch:grpc}).
  \item \textbf{[senior] What is backpressure and where does it appear in
    APIs?}\\
    A fast producer overwhelming a slow consumer. It appears as queue
    growth: in connection pools, in async streams (SSE/WebSockets), and in
    broker consumer lag. Controls: bounded queues, flow-control windows,
    429/503 load shedding, and rate limiting upstream
    (Chapter~\ref{ch:async_apis}).
  \item \textbf{[senior] Walk through a webhook delivery system's failure
    modes.}\\
    Delivery is at-least-once, so consumers must be idempotent; ordering
    is not guaranteed; the consumer's endpoint can be down (retry
    schedule with backoff, dead-letter after $n$ attempts); payloads can
    be spoofed (sign with HMAC, verify before parsing); and slow consumer
    responses are backpressure on your fleet.
  \item \textbf{[staff] How would you design an API platform's governance
    so teams don't route around it?}\\
    Make the compliant path the fastest path: templates that pass linting
    by default, a catalog that gives discoverability in exchange for
    metadata, CI gates that are fast and actionable, and an explicit,
    time-boxed exception process. Governance that only says ``no'' gets
    bypassed; governance that ships value gets adopted
    (Chapter~\ref{ch:lifecycle}).
  \item \textbf{[staff] A partner integrates against your API and you must
    change a field's meaning. Walk the full process.}\\
    You do not change meaning in place. Add the new field alongside the
    old; publish the deprecation with a sunset date; measure usage of the
    old field per consumer; contact stragglers through the developer
    portal; dual-write/dual-read during transition; remove only after
    sunset with usage at zero. Contract thinking turns a breaking change
    into a managed migration (Chapters \ref{ch:versioning},
    \ref{ch:lifecycle}).
  \item \textbf{[staff] Your CEO asks for ``five-nines.'' What do you say?}\\
    Five-nines is 5.26 minutes of downtime \emph{per year}. Price it:
    multi-region active-active, chaos-tested failover, mature on-call,
    and accepting that any single dependency with a weaker SLO caps you.
    Show the marginal cost curve from 99.9 to 99.95 to 99.99 and let the
    business buy the point it needs --- per user journey, not uniformly.
  \item \textbf{[staff] How do you sunset a widely used public API with
    minimal ecosystem damage?}\\
    Announce early with a date and migration guide; sunset headers on
    every response; staged brownouts (scheduled short outages) to flush
    out silent consumers; direct outreach keyed off logs; a final
    hard-stop date you actually honor. The goal is zero surprised
    consumers, because surprised consumers write blog posts.
\end{enumerate}
%% ============================================================================
\section{External Bibliography \& Further Reading}
\label{sec:ch0-reading}

\subsection{Foundational Books}
\begin{itemize}
  \item Roy T. Fielding, \emph{Architectural Styles and the Design of
    Network-based Software Architectures} (2000) \cite{fielding2000rest} ---
    the REST dissertation; read Chapter 5 before Chapter 3 of this book.
  \item Martin Kleppmann, \emph{Designing Data-Intensive Applications}
    \cite{kleppmann2017ddia} --- the distributed-systems foundation beneath
    every failure mode in our taxonomy.
  \item Betsy Beyer et al., \emph{Site Reliability Engineering}
  \cite{beyer2016sre} --- SLOs, error budgets, and the
    operational culture behind Module 4.
  \item Michael T. Nygard, \emph{Release It!}, 2nd ed.
    \cite{nygard2018releaseit} --- stability patterns and the canonical
    timeout/bulkhead/circuit-breaker case studies.
  \item Sam Newman, \emph{Building Microservices}, 2nd ed.
    \cite{newman2021microservices} --- service decomposition and
    inter-service communication for Chapters 11--15.
  \item Lukas Gehl, \emph{API Design Patterns} \cite{gehl2021apidesign} ---
    product-thinking and pattern language for API designers.
\end{itemize}

\subsection{Standards and Specifications}
\begin{itemize}
  \item RFC 9110, \emph{HTTP Semantics} \cite{rfc9110} --- the normative
    reference for methods, status codes, and header fields.
  \item RFC 9111, \emph{HTTP Caching}; RFC 9112, \emph{HTTP/1.1} ---
    companions to RFC 9110 for Chapters 1 and 10.
  \item RFC 9457, \emph{Problem Details for HTTP APIs} --- the error-format
    standard implemented from Chapter 7 onward.
  \item OpenAPI Specification 3.1 \cite{openapi31} --- the contract language
    used from Chapter 4 onward.
  \item The Twelve-Factor App \cite{twelvefactor} --- configuration and
    operational principles assumed by our deployment chapters.
  \item OWASP API Security Top 10 (2023) --- the threat catalog of
    Chapter 16.
\end{itemize}

\subsection{Online Documentation and Courses}
\begin{itemize}
  \item FastAPI official documentation and tutorial --- the framework of
    Module 2.
  \item Python \texttt{http.server}, \texttt{urllib}, and \texttt{logging}
    documentation --- the stdlib surfaces used in this chapter's listings.
  \item Google SRE Workbook (online) --- worked SLO and alerting examples.
  \item MDN Web Docs, \emph{HTTP} reference --- quick, accurate header and
    status-code lookups.
  \item \emph{Designing APIs} (public university API-economy courses and
    conference tracks such as API City and apidays) --- current practice in
    API product management and governance.
\end{itemize}

\subsection{Recommended Datasets (Synthetic Only)}
\begin{itemize}
  \item \textbf{Northwind Robotics catalog} (\texttt{datasets\textbackslash}):
    parts, robots, and telemetry fixtures used throughout the labs;
    entirely fictional, no real PII.
  \item \textbf{Generated latency traces}: produce your own with the
    benchmark harness (Listing~\ref{lst:ch0-bench}) --- real measurements
    from your machine are the best dataset for evaluation exercises.
  \item \textbf{Synthetic event streams} for Chapter 11: generated with
    seeded randomness per that chapter's fixtures; never scraped or real
    customer data.
\end{itemize}

%% ============================================================================
\section{Capstone Projects}
\label{sec:ch0-capstones}

Five projects, one per difficulty band. Each is offline, stdlib-first, and
built on the Northwind Robotics universe. These are portfolio pieces: keep
them in git with measured results committed alongside the code.

\subsection{Project 0.1 --- Contract Archaeology \textnormal{[junior]}}

\textbf{Objective.} Reverse-engineer the complete contract of
Listing~\ref{lst:ch0-miniapi} by observation alone, then write it down
formally.

\textbf{Inputs/Datasets.} The running mini API; the Northwind parts fixture
embedded in it.

\textbf{Steps.}
\begin{enumerate}
  \item Probe every path and method combination; record status codes,
        headers, and bodies.
  \item Draft an informal OpenAPI-style description: paths, parameters,
        response schemas, error shapes.
  \item Deliberately trigger each failure mode and record its body.
  \item Diff your contract against the source; note everything observation
        could not reveal.
\end{enumerate}

\textbf{Acceptance Criteria.}
\begin{itemize}
  \item[$\square$] Written contract covers 100\% of paths and error bodies.
  \item[$\square$] Each operation's failure semantics are enumerated.
  \item[$\square$] A one-paragraph reflection names what black-box probing
    cannot discover (e.g.\ retryability, rate limits).
\end{itemize}

\textbf{Stretch Goals.} Hand-write a valid OpenAPI 3.1 document for the API
and validate it with a local linter.

\subsection{Project 0.2 --- Honest Benchmarks \textnormal{[mid]}}

\textbf{Objective.} Produce a defensible latency study of the mini API under
three load shapes, exposing where averages and percentiles diverge.

\textbf{Inputs/Datasets.} Listings \ref{lst:ch0-bench} and
\ref{lst:ch0-miniapi}; self-generated latency traces.

\textbf{Steps.}
\begin{enumerate}
  \item Extend the harness to support concurrency via threads (keep the
        protocol: warmup, $n \geq 200$, raw samples).
  \item Run three configurations: sequential, 4 concurrent clients, 16
        concurrent clients. Three repetitions each.
  \item Report mean, p50, p95, p99, max, throughput, and error counts per
        configuration in the addendum's benchmark-table format.
  \item Quantify the divergence between mean and p99 as load rises.
\end{enumerate}

\textbf{Acceptance Criteria.}
\begin{itemize}
  \item[$\square$] Nine runs (3 configs $\times$ 3 reps) with raw samples
    retained in CSV.
  \item[$\square$] Written analysis: at what concurrency does p99 detach
    from p50, and which budget stage (Equation~\eqref{eq:ch0-budget}) is
    responsible?
  \item[$\square$] One paragraph on why the mean understated user pain.
\end{itemize}

\textbf{Stretch Goals.} Add payload-size as a variable and fit a simple
model $L = a + b \cdot \text{bytes}$.

\subsection{Project 0.3 --- Failure Taxonomy Field Guide \textnormal{[mid]}}

\textbf{Objective.} Build a runnable gallery: one reproducible, scripted
demonstration per taxonomy category, with detection evidence for each.

\textbf{Inputs/Datasets.} The mini API; a second stdlib service you write to
play ``flaky dependency.''

\textbf{Steps.}
\begin{enumerate}
  \item Script five demos: transport (connection refused / hang), contract
    (schema change), semantic (wrong-unit response), operational (injected
    latency under load), security (missing authz check on an object).
  \item For each, capture the client-visible symptom and the telemetry that
    detects it (status codes, percentiles, logs).
  \item Write the playbook entry an on-call engineer would follow.
\end{enumerate}

\textbf{Acceptance Criteria.}
\begin{itemize}
  \item[$\square$] Five scripts, each deterministic and offline.
  \item[$\square$] Each demo maps symptom $\to$ category $\to$ detection
    strategy per Table~\ref{tab:ch0-taxonomy}.
  \item[$\square$] Playbook entries follow Section~\ref{sec:ch0-playbook}.
\end{itemize}

\textbf{Stretch Goals.} Package the flaky dependency as a reusable
``chaos stub'' with configurable failure probabilities (seeded).

\subsection{Project 0.4 --- Budgets and Breakers for the Mini API
  \textnormal{[senior]}}

\textbf{Objective.} Turn the mini API plus a slow dependency into a
resilience demonstrator: timeouts, retry policy, circuit breaker, and an
SLO with a live error-budget burn calculation.

\textbf{Inputs/Datasets.} Mini API; your chaos stub from Project 0.3 (or a
fresh one with configurable delay/failure injection).

\textbf{Steps.}
\begin{enumerate}
  \item Add a \texttt{/v1/price-check} endpoint that calls the dependency.
  \item Implement client-side: 500\,ms connect / 2\,s read timeouts,
        3-attempt backoff with jitter, and a closed/open/half-open breaker.
  \item Define an SLI (non-5xx within 300\,ms) and a 99.9\% SLO; compute the
        budget from Equation~\eqref{eq:ch0-errorbudget}.
  \item Run a 10-minute scripted failure scenario; plot or tabulate budget
        burn over time and show the breaker's effect.
\end{enumerate}

\textbf{Acceptance Criteria.}
\begin{itemize}
  \item[$\square$] Identical fault scenario run with and without protections;
    the protected run shows materially lower p99 and error rate.
  \item[$\square$] Error-budget computation shown from raw samples.
  \item[$\square$] A post-incident report in the war-story format of
    Section~\ref{sec:ch0-warstory}.
\end{itemize}

\textbf{Stretch Goals.} Add a bulkhead (separate executor for the
dependency call) and demonstrate isolation of an unrelated endpoint.

\subsection{Project 0.5 --- API Platform Governance Starter
  \textnormal{[staff]}}

\textbf{Objective.} Design (and partially automate) the minimum governance
plane that would hold a 20-team organization at maturity L2: catalogs,
contract gates, ownership, and deprecation policy.

\textbf{Inputs/Datasets.} Three or more of your own prior lab APIs as the
``fleet''; the governance baseline of Section~\ref{sec:ch0-governance}.

\textbf{Steps.}
\begin{enumerate}
  \item Define the metadata schema for the catalog (owner, version, lifecycle
        state, contract URI, SLO).
  \item Write a linter (stdlib Python is fine) that scores each API's
        contract document against your rules: has version, has problem
        details, has pagination caps, declares auth scheme.
  \item Author the deprecation policy: timelines, consumer notification,
        sunset mechanics.
  \item Produce a one-page ``maturity report'' per API with L-level and
        next actions.
\end{enumerate}

\textbf{Acceptance Criteria.}
\begin{itemize}
  \item[$\square$] Catalog holds every lab API with complete metadata.
  \item[$\square$] Linter runs offline and exits non-zero on violations
    (CI-ready).
  \item[$\square$] Deprecation policy includes measurable sunset criteria.
  \item[$\square$] Maturity reports agree with a manual audit on at least
    one API.
\end{itemize}

\textbf{Stretch Goals.} Generate a static HTML developer portal from the
catalog (no frameworks; template strings are fine).

%% ============================================================================
\section{Python Dependencies Appendix}
\label{sec:ch0-deps}

This chapter's own listings are stdlib-only. The file below is the baseline
\texttt{requirements.txt} for the whole book's companion environment; later
chapters extend, never replace, it. Pin exactly as shown; the sanity check
(Listing~\ref{lst:ch0-sanity}) verifies importability, not versions, so
minor-version flexibility is acceptable in practice.

\begin{minted}{text}
# requirements.txt -- baseline environment for "Mastering APIs"
# Python 3.11+ ; install with:  pip install -r requirements.txt
fastapi==0.115.6
uvicorn==0.34.0
httpx==0.28.1
pydantic==2.10.4
pytest==8.3.4
\end{minted}

\subsection{fastapi}

\textbf{Description.} The ASGI web framework used from
Chapter~\ref{ch:fastapi} onward for all provider-side services.

\textbf{Why included.} Type-driven request validation, automatic OpenAPI 3.1
generation \cite{openapi31}, and first-class async support make it the
cleanest vehicle for teaching contract-first API design; its generated
schemas plug directly into the contract-testing toolchain of
Chapter~\ref{ch:testing}.

\textbf{Alternatives.} Flask (simpler, but validation and OpenAPI are
add-ons), Django REST Framework (heavier, ORM-coupled), Litestar (comparable,
smaller ecosystem).

\textbf{Installation.}
\begin{minted}{text}
pip install fastapi==0.115.6
\end{minted}

\subsection{uvicorn}

\textbf{Description.} A minimal ASGI server that runs FastAPI applications.

\textbf{Why included.} The reference server for every lab: one process, one
command, deterministic startup on \texttt{127.0.0.1}. Chapter 17 replaces it
behind a production gateway, where its single-process model is discussed
honestly.

\textbf{Alternatives.} Hypercorn (HTTP/2 support), Gunicorn with Uvicorn
workers (POSIX-oriented process management), Daphne.

\textbf{Installation.}
\begin{minted}{text}
pip install uvicorn==0.34.0
\end{minted}

\subsection{httpx}

\textbf{Description.} A typed HTTP client (sync and async) with connection
pooling, HTTP/2 support, and per-request timeouts.

\textbf{Why included.} It is the consumer-side workhorse from
Chapter~\ref{ch:consuming_python}: explicit timeout configuration, transport
plug-ins for offline stubbing, and an API surface that mirrors the protocol
concepts taught in Chapter~\ref{ch:http_wire} rather than hiding them.

\textbf{Alternatives.} \texttt{urllib} (stdlib, used in this chapter to show
the primitives), \texttt{requests} (synchronous only, less explicit about
timeouts), \texttt{aiohttp} (async-first).

\textbf{Installation.}
\begin{minted}{text}
pip install httpx==0.28.1
\end{minted}

\subsection{pydantic}

\textbf{Description.} Data-validation library that parses untrusted input
into typed Python models using standard type annotations.

\textbf{Why included.} It is the concrete implementation of
Chapter~\ref{ch:serialization}'s core idea --- validation at the system
boundary with typed models --- and it powers FastAPI's request/response
contracts. Its error output is also the raw material for RFC~9457 responses
in Chapter~\ref{ch:problem_details}.

\textbf{Alternatives.} \texttt{dataclasses} + manual validation (shown in
this chapter for contrast), \texttt{marshmallow}, \texttt{msgspec}
(faster, narrower).

\textbf{Installation.}
\begin{minted}{text}
pip install pydantic==2.10.4
\end{minted}

\subsection{pytest}

\textbf{Description.} The test runner used by every chapter's test suite.

\textbf{Why included.} Deterministic, fixture-based tests are the book's
acceptance mechanism; labs end with \texttt{pytest -q}, and
Chapter~\ref{ch:testing} builds the full quality-engineering stack
(contract, property-based, load) on top of it.

\textbf{Alternatives.} \texttt{unittest} (stdlib; used in early offline
drills where zero dependencies matter), \texttt{nose2} (unmaintained ---
avoid).

\textbf{Installation.}
\begin{minted}{text}
pip install pytest==8.3.4
\end{minted}

%% ============================================================================
\section{Chapter Summary \& Key Takeaways}
\label{sec:ch0-summary}

\begin{itemize}
  \item An API is a \textbf{versioned contract}: operations,
    representations, and failure semantics over a transport
    (Definition~\ref{def:api}). Hold the product, contract, and
    failure-surface mental models simultaneously.
  \item Incidents classify into five categories --- transport, contract,
    semantic, operational, security --- each with its own detection
    strategy (Table~\ref{tab:ch0-taxonomy}). Classify before you mitigate.
  \item Maturity climbs L0$\to$L4: scripts, documented REST, contract-tested
    CI, resilient observable services, governed platform. Know where your
    systems sit.
  \item Latency is reported in percentiles under a stated load, never as a
    bare average; budgets decompose per Equation~\eqref{eq:ch0-budget}, and
    error budgets follow from SLOs (99.9\% monthly $=$ 43.2 minutes).
  \item Every network call carries an explicit timeout; the war story of
    Section~\ref{sec:ch0-warstory} is what happens when one does not.
  \item The environment is offline-first, deterministic, typed, and
    PowerShell-native; the sanity check gates every lab.
  \item You now have a working stdlib API, client, and benchmark harness ---
    the seed from which twenty-one chapters grow.
\end{itemize}

%% ============================================================================
\section{Quick-Reference Card}
\label{sec:ch0-quickref}

\begin{table}[h]
  \centering
  \small
  \begin{tabularx}{\textwidth}{l X}
    \toprule
    \textbf{Item} & \textbf{Value / rule} \\
    \midrule
    API definition & Versioned contract: operations + representations +
      failure semantics over a transport. \\
    Failure taxonomy & Transport / contract / semantic / operational /
      security --- classify, then mitigate. \\
    Maturity levels & L0 scripts; L1 documented; L2 contract-tested CI;
      L3 resilient + observable; L4 governed platform. \\
    Latency reporting & p50 / p95 / p99, $n \geq 200$, warmup 50, state the
      load; never averages alone. \\
    Latency budget & $L = L_{TLS} + L_{queue} + L_{validate} + L_{logic} +
      L_{serialize} + L_{egress}$. \\
    Error budget & $1 - \mathrm{SLO}$; 99.9\%/month $=$ 43.2 min;
      99.95\% $=$ 21.6 min; 99.99\% $=$ 4.32 min. \\
    Timeouts & Explicit connect + read timeout on every call; no defaults
      trusted. \\
    Retries & Bounded (3), exponential backoff + jitter, idempotent ops
      only, honor \texttt{Retry-After}. \\
    Little's Law & $C = \lambda L$: concurrency $=$ rate $\times$ latency. \\
    STRIDE & Spoofing, Tampering, Repudiation, Information disclosure,
      Denial of service, Elevation of privilege. \\
    Environment & \texttt{python -m venv .venv} then
      \texttt{.\textbackslash .venv\textbackslash Scripts\textbackslash
      Activate.ps1}; run the sanity check before any lab. \\
    Golden rules & Offline-first; deterministic; stdlib-first; synthetic
      data only; PowerShell-native docs. \\
    \bottomrule
  \end{tabularx}
  \caption{Chapter 0 quick-reference card.}
  \label{tab:ch0-quickref}
\end{table}

\section{Standardized Evidence Addendum}
\subsection{Benchmark Snapshot Template}
\begin{table}[h]
  \centering
  \small
  \begin{tabularx}{\textwidth}{l c c c c}
    \toprule
    \textbf{Config} & \textbf{p50 latency} & \textbf{p99 latency} & \textbf{Error rate} & \textbf{Throughput} \\
    \midrule
    Baseline & -- & -- & -- & 1.00x \\
    Candidate A & -- & -- & -- & -- \\
    Candidate B & -- & -- & -- & -- \\
    \bottomrule
  \end{tabularx}
  \caption{Minimum benchmark table to keep per chapter.}
\end{table}

\subsection{Request Trace Template}
\begin{itemize}
  \item \textbf{Request:} one representative API call for this chapter's topic.
  \item \textbf{Before:} behavior (headers, payload, latency, failure mode) with the naive approach.
  \item \textbf{After:} behavior with the technique in this chapter.
  \item \textbf{Result:} what changed in correctness, resilience, latency, and operability.
\end{itemize}

\subsection{Cost and Latency Budget Snapshot}
\begin{table}[h]
  \centering
  \small
  \begin{tabularx}{\textwidth}{l c c}
    \toprule
    \textbf{Stage} & \textbf{Latency budget} & \textbf{Cost driver} \\
    \midrule
    TLS / connection setup & -- & handshake round trips, keep-alive hit rate \\
    Request validation & -- & schema complexity, CPU per request \\
    Business logic and I/O & -- & downstream fan-out, query cost \\
    Serialization and egress & -- & payload size, compression ratio \\
    \bottomrule
  \end{tabularx}
  \caption{Operational budget template for this chapter's technique.}
\end{table}

\section{Configuration Preset Template}
\begin{minted}{text}
# api_policy.yaml
policy_version: "v1.0.0"
component: "<chapter_topic>"
timeout_ms: 0
retry_budget: 0
fallback_strategy: "fail_fast"
rollback_trigger: "error_budget_burn_gt_threshold"
\end{minted}

\section{CI Quality Gate Specification}
\begin{itemize}
  \item contract tests green against the pinned OpenAPI/AsyncAPI/Protobuf spec,
  \item no breaking schema changes without a version bump,
  \item error responses conform to RFC 9457 problem-details shape,
  \item benchmark deltas within approved regression bounds (latency and throughput),
  \item policy version and rollback plan present in release metadata.
\end{itemize}

\section{Incident Runbook Template}
\begin{enumerate}
  \item Detect regression from SLO burn-rate alerts or consumer reports.
  \item Scope impacted endpoints, versions, and consumer segments.
  \item Contain impact (rollback, feature flag, rate-limit tightening, or traffic shift).
  \item Diagnose with traces, structured logs, and contract diffs.
  \item Recover with validated fix and verified rollout procedure.
  \item Prevent recurrence via new regression tests and CI gates.
\end{enumerate}

\section{Cross-Chapter Decision Card}
\begin{table}[h]
  \centering
  \small
  \begin{tabularx}{\textwidth}{l X X}
    \toprule
    \textbf{Dimension} & \textbf{Choose this approach when} & \textbf{Prefer an alternative when} \\
    \midrule
    Use case & -- & -- \\
    Latency budget & -- & -- \\
    Operational cost & -- & -- \\
    Team maturity & -- & -- \\
    \bottomrule
  \end{tabularx}
  \caption{Decision card to complete per chapter.}
\end{table}

\section{ROI Model and Build-vs-Buy}
\begin{itemize}
  \item \textbf{Build when:} the capability is differentiating, compliance forbids vendors, or scale makes vendor pricing irrational.
  \item \textbf{Buy when:} the capability is commodity (auth, gateways, observability), time-to-market dominates, or staffing is thin.
  \item \textbf{Hidden costs to model:} on-call load, upgrade cadence, expertise ramp, data egress fees.
\end{itemize}

\section{Disaster Recovery Notes (RPO/RTO)}
\begin{itemize}
  \item Define recovery point objective (RPO): maximum acceptable data loss window.
  \item Define recovery time objective (RTO): maximum acceptable restoration time.
  \item Test restores regularly; an untested backup is not a backup.
\end{itemize}

